\documentclass[11pt,a4paper]{article}

\usepackage[T1]{fontenc}
\usepackage[utf8]{inputenc}
\usepackage{amsmath,amssymb,amsthm,mathrsfs}
\usepackage{bm}
\usepackage[margin=1in]{geometry}
\usepackage{hyperref}
\usepackage{cite}
\usepackage{enumitem}

\hypersetup{colorlinks=true,linkcolor=blue,citecolor=blue,urlcolor=blue}

\theoremstyle{plain}
\newtheorem{theorem}{Theorem}[section]
\newtheorem{proposition}[theorem]{Proposition}
\newtheorem{lemma}[theorem]{Lemma}
\newtheorem{corollary}[theorem]{Corollary}
\theoremstyle{definition}
\newtheorem{definition}[theorem]{Definition}
\newtheorem{assumption}[theorem]{Assumption}
\newtheorem{remark}[theorem]{Remark}
\newtheorem{example}[theorem]{Example}

\newcommand{\kfive}{\kappa_5}
\newcommand{\kfour}{\kappa_4}
\newcommand{\Lfive}{\Lambda_5}
\newcommand{\Lfour}{\Lambda_4}
\newcommand{\TT}{\mathrm{TT}}
\newcommand{\KK}{\mathrm{KK}}
\newcommand{\RS}{\mathrm{RS}}
\newcommand{\Diff}{\mathrm{Diff}}
\newcommand{\Tr}{\operatorname{Tr}}
\newcommand{\Op}{\operatorname{Op}}
\newcommand{\ee}{\mathbf{e}}

\title{\bf Bulk-to-Brane Projection, Dynamical Nonclosure, and Observable
Residues in Randall--Sundrum Gravity}

\author{%
  Jeremy Rodgers\thanks{Website: \href{https://everythingequation.com}{everythingequation.com}}\\
  \small Independent Researcher
}

\date{June 2026}

\begin{document}

\maketitle

\begin{abstract}
We study, within the positive-tension, two-sided, $Z_2$-symmetric
Randall--Sundrum model (RS2), the precise sense in which the
instantaneous state of a brane observer fails to determine its own
future, and the precise sense in which that failure can and cannot be
detected. Selecting RS2 is an additional physical postulate; nothing
below derives it. For a brane of tension $\lambda>0$ with
five-dimensional coupling $\kfive^2$, bulk cosmological constant
$\Lfive$, induced metric $q_{\mu\nu}$, and brane stress tensor
$T_{\mu\nu}$, the Israel condition and the Gauss--Codazzi identity give
the exact projected equation
$G_{\mu\nu}[q]+\Lfour q_{\mu\nu}=\kfour^2T_{\mu\nu}
+\mathcal C^{\RS}_{\mu\nu}$ with
$\mathcal C^{\RS}_{\mu\nu}=\kfive^4\Pi_{\mu\nu}-E_{\mu\nu}$,
where $\Pi_{\mu\nu}$ is quadratic in $T_{\mu\nu}$ and $E_{\mu\nu}$ is
the electric projection of the bulk Weyl tensor. We construct an
explicit gauge-inequivalent pair of bulk data with identical brane
readout and different brane futures, proving that the source evolution
admits no deterministic quotient flow on the instantaneous brane Cauchy
state over the full RS2 source domain. We then evaluate the projection
exactly in the homogeneous Friedmann sector, obtaining a positive
$\rho^2$ term and a bulk-state-dependent dark-radiation term $\mu/a^4$.
In the linear weak-field sector the Gaussian/Schur-complement elimination
of the bulk is exact at linear tree level; evaluating the resulting
response asymptotically for $r\gg\ell$ --- a controlled expansion, not an
exact closed form at arbitrary distance --- gives the long-distance
potentials
$\Phi=-\tfrac{G_4M}{r}(1+\tfrac{2\ell^2}{3r^2}+\cdots)$ and
$\Psi=-\tfrac{G_4M}{r}(1+\tfrac{\ell^2}{3r^2}+\cdots)$.
Eliminating the common AdS radius $\ell$ between the two regimes yields
the parameter-free cross-regime relation
$A_\Psi=\tfrac12A_\Phi$, $B=\tfrac{\kfour^4}{24}A_\Phi
=\tfrac{\kfour^4}{12}A_\Psi$,
where $H^2\supset B\rho^2$ and $A_\Phi,A_\Psi$ are the leading $1/r^3$
potential coefficients. Violation falsifies minimal RS2; agreement at a
nonzero value is a parameter-saving cross-regime success but does not
falsify a four-dimensional theory that treats the three coefficients as
independent. Finally, an elementary pushforward lemma fixes the exact
interpretive boundary: for every brane-only protocol, including adaptive
ones, the induced Tier-1 history law reproduces the RS2 record
distribution exactly, and in RS2 this
obstruction is realized concretely by the four-dimensional
gravity-plus-CFT description, which reproduces the same $1/r^3$
residue. The result is therefore positive at the level of model-relative
projection and nonclosure, and negative at the level of unrestricted
ontological discrimination.
\end{abstract}

\tableofcontents

\section{Introduction}
\label{sec:intro}

\subsection{The physical problem}

A brane observer in a braneworld model has direct access to the induced
metric on the brane and to the matter fields confined there. Gravity,
however, propagates in the bulk. Two questions immediately arise, and
they are logically independent.

The first is a question about \emph{states}. Does the instantaneous
brane data --- induced three-metric, its brane-intrinsic extrinsic
curvature, and matter Cauchy data --- determine the subsequent brane
history? If not, the brane state is incomplete relative to brane
evolution: there exist bulk configurations, physically distinct after
every gauge identification has been made, which look identical now and
diverge later.

The second is a question about \emph{theories}. Suppose the answer to
the first question is negative. Does it follow that a brane observer
could, even in principle, discover that the world is five-dimensional?
This does not follow, and the distinction is the central logical point
of this paper. Failure of a deterministic flow on the \emph{instantaneous}
state is not failure of every brane-only history law. Any source theory,
pushed forward through any readout map, defines a brane-only stochastic
process with exactly the correct statistics; that process is generally
nonlocal, non-Markovian, and state-dependent, but it is a complete
theory of brane records.

We answer both questions inside the specific model that makes them
sharpest.

\subsection{Contributions}

\begin{enumerate}[leftmargin=*]
\item \textbf{An explicit same-readout witness.} We exhibit two
linearized RS2 source data sets which agree exactly on the instantaneous
brane readout, are \emph{not} related by any five-dimensional
diffeomorphism (their gauge-invariant Weyl curvature differs), and whose
brane futures differ. Existence of such data is proved, not asserted, by
an explicit Kaluza--Klein spectral construction using the positivity of
the brane spectral weight $u_m(0)^2$. This establishes failure of
deterministic projectability of the source flow onto the instantaneous
brane Cauchy state, over the full RS2 source domain
(Section~\ref{sec:witness}, Theorem~\ref{thm:main} part 1). The result is conditional on that source
domain: if the model is \emph{defined} to admit only the stationary
vacuum bulk state, the witness lies outside the domain
(Section~\ref{sec:failure}, item 1).

\item \textbf{Derivation, not quotation, of the projected equation.}
Section~\ref{sec:projected} carries out the Israel matching and
Gauss--Codazzi projection in full, including all coefficients, signs, and
dimensions, and derives the Bianchi/nonclosure relation
$\nabla^\mu E_{\mu\nu}=\kfive^4\nabla^\mu\Pi_{\mu\nu}$, which constrains
the divergence of $E_{\mu\nu}$ but not its transverse-traceless content.

\item \textbf{Complete evaluation in two regimes.} The homogeneous sector
(Section~\ref{sec:frw}) is evaluated exactly. In the linear weak-field
sector (Section~\ref{sec:linear}) two distinct claims must be kept apart:
the Gaussian/Schur-complement elimination of the bulk is \emph{exact at
linear tree level}, whereas the resulting metric potentials
\eqref{eq:potentials} are a \emph{controlled asymptotic evaluation} valid
for $r\gg\ell$ and not an exact expression at arbitrary distance. Brane
bending is treated explicitly throughout.

\item \textbf{A parameter-eliminated cross-regime relation.} Eliminating
$\ell$ between the two regimes gives
$A_\Psi=A_\Phi/2$ and $B=\kfour^4A_\Phi/24$
(Section~\ref{sec:crossregime}). We claim no literature priority for this
algebraic elimination; the two input results are established. Its content
is a falsifiable, parameter-free link between two otherwise independent
observational sectors.

\item \textbf{An exact operational-discrimination boundary.}
Section~\ref{sec:opeq} gives the classical pushforward lemma --- whose
proof is immediate, and whose value is interpretive rather than technical
--- together with its quantum process-tensor and closed-time-path forms,
and states exactly what is and is not excluded (Section~\ref{sec:main}).
\end{enumerate}

\subsection{Relation to established results, and this paper's contribution}
\label{sec:relation}

The physical inputs used below are established results and are credited as
such. The RS2 model and its tuning are due to Randall and Sundrum
\cite{RS2}. The junction condition is Israel's \cite{Israel}, and the
projected brane equation with $\Pi_{\mu\nu}$ and $E_{\mu\nu}$ is that of
Shiromizu, Maeda, and Sasaki \cite{SMS}. The brane Friedmann equation with
its $\rho^2$ and $\mu/a^4$ terms is due to Binetruy, Deffayet, Ellwanger,
and Langlois \cite{BDEL} and Mukohyama \cite{Mukohyama1,Mukohyama2}. The
linearized brane propagator, the KK spectral weight, the brane-bending
treatment, and the long-distance potentials are due to Garriga and Tanaka
\cite{GarrigaTanaka} and Giddings, Katz, and Randall \cite{GKR}. The
four-dimensional CFT reproduction of the $1/r^3$ residue is due to Duff
and Liu \cite{DuffLiu} and related work
\cite{APR,PerezVictoria}. The influence-functional
\cite{FeynmanVernon}, Mori--Zwanzig \cite{MoriZwanzig}, and process-tensor
\cite{Pollock} formalisms are standard.

The distinctive contribution of this paper is their integration into a
single proved chain. Concretely, we establish that
\begin{equation}
\label{eq:chain}
\text{same instantaneous brane state}
\;\not\Rightarrow\;
\text{same brane future},
\end{equation}
by way of
\begin{equation}
\label{eq:chain2}
\begin{gathered}
\text{explicit gauge-reduced bulk witness}
\;+\;
\mathcal C^{\RS}_{\mu\nu}=\kfive^4\Pi_{\mu\nu}-E_{\mu\nu}\\
+\;\text{parameter-eliminated cross-regime relation}
\;+\;\text{exact operational-discrimination boundary}.
\end{gathered}
\end{equation}
The projected equation and the weak-field and cosmological inputs are
established physics; what is new here is the rigorous combined
source/readout argument \cite{Paper1,Paper2,Paper3,Paper4,Paper5,Paper6},
the explicit spectral construction of the
same-readout witness, the cross-regime algebraic synthesis, and the exact
statement of its interpretive boundary. We make no literature-priority
claim for the coefficient elimination itself.

\subsection{Scope, and what is \emph{not} claimed}
\label{sec:scope}

RS2 is not derived from any prior structure; its selection is a physical
postulate. Nothing below shows that our universe \emph{is} a projection.
The cross-regime relation does not prove five-dimensional ontology:
Section~\ref{sec:fourdim} exhibits a four-dimensional realization of the
same weak-field residue, and Section~\ref{sec:opeq} proves that no
brane-only experiment can do better. The KK branch cut excludes only a
pre-specified finite-pole free or tree-level linear alternative, not every
local interacting four-dimensional theory. The long-distance weak-field
expansion is not valid at $r\lesssim\ell$, and the bulk parameter $\mu$ is
not fixed without a bulk-state condition. All nine conditions on which the
positive results depend --- including the bulk-state domain, the RS
tuning, positive tension, the no-flux assumption, and the detuned-brane
limitation --- are collected once, in Section~\ref{sec:failure}, and are
not repeated elsewhere.

\section{The five-dimensional source model}
\label{sec:model}

\subsection{Action and conventions}

Let $(\mathcal M_5,g_{AB})$ be a five-dimensional Lorentzian manifold
containing a timelike hypersurface $\Sigma$ (the brane). Capital Latin
indices $A,B,\dots$ run over the five dimensions; Greek indices
$\mu,\nu,\dots$ are tangent to $\Sigma$; Latin $i,j,\dots$ are spatial on
a brane Cauchy slice. The signature is $(-++++)$, and $c=\hbar=1$. The
classical action is
\begin{equation}
\label{eq:action}
\begin{aligned}
S_5[g,X,\psi]
={}&\frac{1}{2\kfive^2}\int_{\mathcal M_5}d^5x\,\sqrt{-g_5}\,
\bigl(R_5-2\Lfive\bigr)\\
&+\frac{1}{\kfive^2}\sum_{\pm}\int_{\Sigma}d^4x\,\sqrt{-q}\,K_\pm
+\int_{\Sigma}d^4x\,\sqrt{-q}\,
\bigl[-\lambda+\mathcal L_m(q,\psi)\bigr].
\end{aligned}
\end{equation}
Gibbons--Hawking--York terms appear on \emph{both} sides of the thin
brane, with the corresponding outward normals. The embedding is
$X^A(x^\mu)$, the tangent basis $e^A{}_\mu=\partial_\mu X^A$, the unit
spacelike normal $n^A$, and
\begin{equation}
\label{eq:induced}
q_{\mu\nu}=g_{AB}e^A{}_\mu e^B{}_\nu,
\qquad
K_{\mu\nu}=e^A{}_\mu e^B{}_\nu\nabla_A n_B .
\end{equation}
The sign convention for $\Lfive$ is fixed by requiring the bulk field
equation away from the brane to read
\begin{equation}
\label{eq:bulkEE}
{}^{(5)}G_{AB}+\Lfive g_{AB}=0 ,
\end{equation}
so $[\Lfive]=L^{-2}$. The brane stress tensor is
\begin{equation}
\label{eq:Tdef}
T_{\mu\nu}=-\frac{2}{\sqrt{-q}}
\frac{\delta(\sqrt{-q}\,\mathcal L_m)}{\delta q^{\mu\nu}} .
\end{equation}
Dimensionally, $[\kfive^2]=L^3$, $[\kfour^2]=L^2$, $[T_{\mu\nu}]=L^{-4}$,
and $[E_{\mu\nu}]=L^{-2}$. The brane tension appears in \eqref{eq:action}
integrated over $d^4x$, exactly like $\mathcal L_m$; it is therefore a
four-dimensional energy density,
\begin{equation}
\label{eq:lambdadim}
[\lambda]=L^{-4},
\qquad
[\kfive^2\lambda]=L^{3}\cdot L^{-4}=L^{-1}.
\end{equation}
This is consistent throughout: $K_{\mu\nu}\sim\kfive^2\lambda$ has
dimension $L^{-1}$, as an extrinsic curvature must; and
$\kfour^2=\kfive^4\lambda/6$ gives
$[\kfour^2]=L^{6}\cdot L^{-4}=L^{2}$, as required.

Varying \eqref{eq:action} with respect to $g_{AB}$ gives
\eqref{eq:bulkEE} in the bulk together with the Israel junction condition
at $\Sigma$; varying with respect to $X^A$ gives the brane equation of
motion, which for a $Z_2$-symmetric configuration is implied by the
junction condition and the Codazzi identity. The response tensor
$\mathcal C^{\RS}_{\mu\nu}$ below is obtained from this variation; it is
not introduced by hand.

\subsection{Physical equivalence and the reduced source space}

\begin{definition}[Source configuration]
\label{def:source}
A source configuration is a tuple
\begin{equation}
\label{eq:sourcetuple}
s=\bigl(g_{AB},\,X^A,\,\psi;\ \text{initial and boundary data}\bigr),
\end{equation}
solving the five-dimensional constraints, with the brane embedded as a
timelike hypersurface.
\end{definition}

\begin{definition}[Physical equivalence]
\label{def:equiv}
Two source configurations are physically equivalent when related by a
five-dimensional diffeomorphism carrying the brane and its orientation to
the same physical hypersurface, composed with an internal gauge
transformation of the brane matter fields. The reduced source space is
\begin{equation}
\label{eq:Sspace}
S_{\RS}
=\frac{\{(g_{AB},X^A,\psi;\ \text{data})\ \text{solving the constraints}\}}
{\Diff_5\ltimes\mathcal G_m}.
\end{equation}
\end{definition}

This quotient does \emph{not} remove bulk transverse-traceless
gravitational waves, Kaluza--Klein continuum amplitudes, or the
AdS--Schwarzschild mass parameter. Each changes gauge-invariant
five-dimensional curvature invariants, energy flux, horizon structure, or
the later induced brane metric, and hence is not pure gauge. This point
is what makes the witness of Section~\ref{sec:witness} a genuine physical
distinction rather than a coordinate artifact.

\subsection{The instantaneous brane readout}
\label{sec:readout}

Choose a five-dimensional Cauchy surface $\mathcal C_5$ intersecting the
brane worldvolume in a three-dimensional surface $\mathcal C_3$.

\begin{definition}[Tier-1 readout]
\label{def:readout}
The instantaneous brane readout is the map
\begin{equation}
\label{eq:readout}
p:S_{\RS}\longrightarrow T_{\mathrm{brane}},
\qquad
p(s)=\bigl(\gamma_{ij},\Theta_{ij};\ \psi,\pi_\psi\bigr)\big|_{\mathcal C_3},
\end{equation}
where $\gamma_{ij}$ is the three-metric induced on $\mathcal C_3$,
$\Theta_{ij}$ is the extrinsic curvature of $\mathcal C_3$ regarded as a
spatial hypersurface \emph{within the brane worldvolume}, and
$(\psi,\pi_\psi)$ are the brane matter Cauchy data. The constraints
induced from the five-dimensional equations and the matter gauge
constraints are imposed; spatial diffeomorphisms of $\mathcal C_3$ and
matter gauge transformations are quotiented.
\end{definition}

Two features of this definition matter. First, $\Theta_{ij}$ records the
gauge-independent first normal derivative of the induced metric along the
brane without presupposing that the brane possesses a closed intrinsic
Einstein Hamiltonian --- which is exactly what is at issue. Second, the
quotient by spatial diffeomorphisms and the imposition of the constraints
means that no lapse-, shift-, or time-coordinate-dependent velocity
enters, so the readout is not made noninjective by a trivial gauge
artifact.

The readout is physically motivated inside the model: standard-model
matter is confined to $\Sigma$, and local rods, clocks, detectors, and
gravitational-wave interferometers respond to induced brane fields. Bulk
data can be inferred only through their subsequent influence on those
fields; they are not components of the present brane state.

\begin{remark}[Instantaneity is essential]
\label{rem:instant}
If the readout were redefined to include the \emph{complete} brane
history for all times, any bulk distinction that changes that history
would by definition already be distinguished. The dynamical question is
precisely whether the present brane state closes under time evolution.
Sections~\ref{sec:witness} and \ref{sec:opeq} answer the two versions of
the question in opposite directions, and it is essential not to conflate
them.
\end{remark}

\section{A genuine same-readout bulk distinction}
\label{sec:witness}

Linearize about the positive-tension RS2 background,
\begin{equation}
\label{eq:background}
ds^2=dy^2+e^{-2|y|/\ell}\eta_{\mu\nu}dx^\mu dx^\nu,
\qquad
\Lfive=-\frac{6}{\ell^2},
\qquad
\lambda=\frac{6}{\kfive^2\ell},
\end{equation}
which is the tuning that makes the brane worldvolume flat. After fixing
five-dimensional diffeomorphisms, the bulk perturbation retains
gauge-invariant transverse-traceless (TT) gravitational modes.

\begin{proposition}[Same readout, different future]
\label{prop:witness}
Let the admissible source domain include smooth, finite-energy TT bulk
initial data. Consider
\begin{align}
s_0&:\quad h^{\TT}_{AB}=0,
\label{eq:s0}\\
s_1&:\quad h^{\TT}_{AB}\ \text{a smooth wave packet whose initial
support lies a positive bulk distance from }\Sigma,
\label{eq:s1}
\end{align}
with the packet chosen so that $h^{\TT}_{AB}$ and its first time
derivative vanish identically in a neighborhood of the brane on the
initial Cauchy surface. Then
\begin{equation}
\label{eq:samereadout}
p(s_0)=p(s_1),
\end{equation}
the configurations $s_0$ and $s_1$ are not physically equivalent in the
sense of Definition~\ref{def:equiv}, and there exists $\tau>0$ with
\begin{equation}
\label{eq:diffFuture}
p(\Phi_\tau s_0)\ne p(\Phi_\tau s_1),
\end{equation}
where $\Phi_t$ is the source evolution flow.
\end{proposition}

\begin{proof}
Equation \eqref{eq:samereadout} holds because the readout
\eqref{eq:readout} depends only on data at $\mathcal C_3\subset\Sigma$
and on a neighborhood of $\mathcal C_3$ within the brane, and by
construction $h^{\TT}$ and $\partial_t h^{\TT}$ vanish there; the induced
three-metric, its brane-intrinsic extrinsic curvature, and the matter
data therefore coincide.

The two configurations are not gauge-related. The linearized
five-dimensional Weyl tensor of $s_1$ is nonzero on the support of the
packet, while that of $s_0$ vanishes identically. Weyl curvature is a
tensor; a diffeomorphism cannot map a configuration with vanishing
gauge-invariant curvature to one with nonvanishing curvature. Hence
$[s_0]\ne[s_1]$ in $S_{\RS}$.

It remains to prove \eqref{eq:diffFuture}, i.e.\ that a packet with the
stated properties \emph{exists} and does eventually disturb the brane.
This is established by an explicit spectral construction, which we now
give.

\emph{(i) The homogeneous brane boundary condition.} Set the brane source
in \eqref{eq:bulkwave} to zero and integrate that equation across the
brane over $y\in(-\epsilon,\epsilon)$. The $\partial_y^2$ term gives
$\partial_yh^{\TT}_{\mu\nu}\big|_{0^-}^{0^+}$; the $\delta(y)$ term gives
$(4/\ell)h^{\TT}_{\mu\nu}(0)$; the remaining terms are bounded and drop
out as $\epsilon\to0$. Using $Z_2$ symmetry,
$\partial_yh^{\TT}\big|_{0^-}=-\partial_yh^{\TT}\big|_{0^+}$, so
\begin{equation}
\label{eq:homobc}
\Bigl(\partial_y+\frac{2}{\ell}\Bigr)h^{\TT}_{\mu\nu}\Big|_{y=0^+}=0 ,
\end{equation}
in the conventions of \eqref{eq:bulkwave}. This is the homogeneous RS2
boundary condition.

\emph{(ii) The mode basis.} Let $\{u_m(y)\}$ be the complete set of
eigenfunctions of the transverse operator appearing in
\eqref{eq:bulkwave}, defined on the physical half-space $y\ge0$ with the
homogeneous Robin condition \eqref{eq:homobc} at $y=0^+$ and the
appropriate regularity condition in the bulk interior, and extended to the
two-sided covering space by the required $Z_2$ reflection. The spectrum
consists of
\begin{itemize}[leftmargin=*]
\item one \emph{normalizable discrete} zero mode
$u_0\propto e^{-2|y|/\ell}$; and
\item a \emph{delta-normalized generalized continuum} $\{u_m\}_{m>0}$ of
Bessel combinations satisfying \eqref{eq:homobc} and the bulk/horizon
condition.
\end{itemize}
The continuum modes are \emph{not} normalizable; they are delta-normalized
in the usual generalized sense, and together with $u_0$ they furnish a
complete resolution of the identity on the relevant space of TT bulk data
\cite{RS2,GarrigaTanaka,GKR}.

\emph{(iii) Time-symmetric continuum data.} Choose
\begin{equation}
\label{eq:tsym}
\partial_t h^{\TT}_{\mu\nu}\big|_{t=0}=0,
\end{equation}
together with a nonzero smooth, finite-energy TT displacement
$h^{\TT}_{\mu\nu}\big|_{t=0}$ whose support lies in $y\ge y_0>0$ and which
is \emph{orthogonal to the normalizable zero mode}, so that it lies
entirely in the KK continuum subspace. Such data exist: the zero-mode
overlap is a single linear functional on the displacement, so given any
two smooth profiles supported in disjoint regions of $y\ge y_0$ with
nonzero and unequal zero-mode overlaps, an appropriate linear combination
of them annihilates that functional while remaining nonzero, smooth, and
supported away from the brane. Because the data vanish identically in a
neighborhood of the brane, all compatibility conditions at $y=0$ hold
trivially and the initial brane readout agrees with that of the zero
solution --- which is \eqref{eq:samereadout}. The vanishing initial
velocity \eqref{eq:tsym} makes the resulting solution, and hence its brane
trace, \emph{even in time}.

\emph{(iv) The evolved brane trace.} Expanding the data of (iii) in the
basis of (ii) and in spatial Fourier modes, and using that the zero-mode
coefficient vanishes by construction while
$\widehat g_{\mu\nu}\equiv0$ by \eqref{eq:tsym}, the evolved solution on
the brane is carried entirely by the continuum:
\begin{equation}
\label{eq:branetrace}
h^{\TT}_{\mu\nu}(t,\mathbf x,0)
=\int d^3k\int_0^\infty dm\;
u_m(0)\,\widehat f_{\mu\nu}(m,\mathbf k)\,
\cos\!\bigl(\sqrt{\mathbf k^2+m^2}\;t\bigr)\,
e^{i\mathbf k\cdot\mathbf x},
\end{equation}
where $\widehat f_{\mu\nu}(m,\mathbf k)$ is the continuum spectral
coefficient of the initial displacement.

\emph{(v) Nonvanishing coupled component.} By \eqref{eq:um},
\begin{equation}
\label{eq:upos}
u_m(0)^2
=\frac{2}{\pi^2m\ell\bigl[J_1^2(m\ell)+Y_1^2(m\ell)\bigr]}>0
\qquad\text{for all }m>0,
\end{equation}
so \emph{no} continuum mode is decoupled from the brane. Moreover, by
completeness of the basis of (ii) and the associated Plancherel identity,
the spectral transform is an isometry onto the direct sum of the discrete
and continuum spectral spaces. Since the displacement chosen in (iii) is
nonzero and orthogonal to the discrete zero mode, its continuum spectral
coefficient $\widehat f_{\mu\nu}$ is therefore nonzero on a set of
positive spectral measure.

\emph{(vi) Nonvanishing in the future.} Suppose, for contradiction, that
$h^{\TT}_{\mu\nu}(t,\mathbf x,0)=0$ for every $t\ge0$ and every
$\mathbf x$. By the evenness in time established in (iii), the brane trace
then vanishes for every $t\in\mathbb R$. Fourier-transforming
\eqref{eq:branetrace} in $\mathbf x$, each $\mathbf k$-component vanishes
for all $t$. For fixed $\mathbf k$ the map
\begin{equation}
\label{eq:momega}
m\longmapsto\omega=\sqrt{\mathbf k^2+m^2}
\end{equation}
is one-to-one on $m\ge0$, so \eqref{eq:branetrace} is, for each
$\mathbf k$, a Fourier-cosine transform in the variable $\omega$ of the
reparameterized coefficient. Injectivity of the Fourier-cosine transform
therefore forces
\begin{equation}
\label{eq:forced}
u_m(0)\,\widehat f_{\mu\nu}(m,\mathbf k)=0
\qquad\text{for almost every }(m,\mathbf k).
\end{equation}
Since $u_m(0)^2>0$ for every $m>0$ by \eqref{eq:upos}, this forces
$\widehat f_{\mu\nu}=0$ almost everywhere, contradicting (v). Hence the
brane trace is nonzero at some $t=\tau>0$. Because the induced brane
perturbation and the electric Weyl term $E_{\mu\nu}$ on the brane are
determined by that trace and its normal derivative, \eqref{eq:diffFuture}
follows.

\emph{(vii) Causal delay.} Finally, the linearized bulk equation is
hyperbolic with finite propagation speed on this background, so the brane
trace remains \emph{zero} until the causal future of the remote initial
support reaches $y=0$. The disturbance is therefore genuinely delayed, not
instantaneous, which is exactly why the two configurations share a readout
now and differ later. Together with (vi) this establishes the required
$\tau>0$ same-readout/different-future statement.
\end{proof}

\begin{corollary}[Nonprojectability]
\label{cor:nonproj}
There exists no map $U_t$ on the instantaneous brane readout satisfying
$p\circ\Phi_t=U_t\circ p$ for all admissible source data. Moreover, for
any admissible preparation class rich enough to resolve the two witness
states --- for instance one containing delta preparations at $s_0$ and
$s_1$, or any class whose members assign them different conditional
future laws --- universal (strong) Markov lumpability of the readout
fails as well. Nothing here asserts that a deterministic witness by
itself furnishes a stochastic experiment; that requires a preparation
law.
\end{corollary}

\begin{proof}
Such a $U_t$ would require $p(s_0)=p(s_1)\Rightarrow
p(\Phi_ts_0)=p(\Phi_ts_1)$, contradicting
Proposition~\ref{prop:witness}. The stochastic statement is conditional
as stated: a deterministic witness does not by itself define a stochastic
experiment. It becomes a failure of lumpability once a preparation law is
specified, and \eqref{eq:noflow} then defeats it because the two
conditional future laws within the fiber differ.
\end{proof}

\begin{remark}[Scope]
\label{rem:witnessScope}
Corollary~\ref{cor:nonproj} is conditional on the source domain admitting
the normalizable bulk radiation used in the construction. If the model is
\emph{defined} to have a unique stationary vacuum bulk state with all
homogeneous incoming data excluded, the witness is outside the domain and
the fiber must be reassessed for that restricted domain; the forced
retarded KK response and the effective-action results of
Sections~\ref{sec:linear}--\ref{sec:opeq} survive unchanged. See
Section~\ref{sec:failure}, item 1.
\end{remark}

The construction also identifies \emph{what} must be added. At linear
order the required completion is the bulk gravitational initial data
modulo the subspace whose retarded brane response vanishes identically;
in a KK decomposition, the coupled continuum amplitudes modulo
zero-response modes. In the nonlinear theory it is the bulk
initial/boundary data or, after elimination, the induced nonlocal brane
history functional. Adding the instantaneous value of $E_{\mu\nu}$ alone
does \emph{not} in general close its future evolution --- as
Section~\ref{sec:conservation} makes precise.

\section{The exact projected Einstein equation}
\label{sec:projected}

\subsection{Israel junction condition}

Let the brane surface stress tensor be
\begin{equation}
\label{eq:Smunu}
S_{\mu\nu}=-\lambda q_{\mu\nu}+T_{\mu\nu},
\qquad
S\equiv q^{\mu\nu}S_{\mu\nu}=-4\lambda+T .
\end{equation}
The Israel condition \cite{Israel} for a thin shell in five dimensions is
\begin{equation}
\label{eq:israel}
[K_{\mu\nu}]\equiv K^+_{\mu\nu}-K^-_{\mu\nu}
=-\kfive^2\Bigl(S_{\mu\nu}-\tfrac13q_{\mu\nu}S\Bigr),
\end{equation}
the factor $\tfrac13$ being $1/(D-2)$ with $D=5$. Imposing the RS2 $Z_2$
reflection symmetry gives $K^+_{\mu\nu}=-K^-_{\mu\nu}$, so
$[K_{\mu\nu}]=2K_{\mu\nu}$ evaluated with the outward normal on one side,
and hence
\begin{equation}
\label{eq:Kfixed}
\boxed{\;
K_{\mu\nu}
=-\frac{\kfive^2}{2}
\Bigl[\,T_{\mu\nu}+\tfrac13(\lambda-T)q_{\mu\nu}\,\Bigr].
\;}
\end{equation}
Indeed, $S_{\mu\nu}-\tfrac13q_{\mu\nu}S
=T_{\mu\nu}-\lambda q_{\mu\nu}
-\tfrac13q_{\mu\nu}(T-4\lambda)
=T_{\mu\nu}+\tfrac13q_{\mu\nu}(\lambda-T)$,
which gives \eqref{eq:Kfixed} on halving. Traces:
\begin{equation}
\label{eq:Ktrace}
K=q^{\mu\nu}K_{\mu\nu}
=-\frac{\kfive^2}{2}\Bigl[T+\tfrac43(\lambda-T)\Bigr]
=-\frac{\kfive^2}{2}\Bigl[\tfrac43\lambda-\tfrac13T\Bigr]
=-\frac{\kfive^2}{6}\bigl(4\lambda-T\bigr).
\end{equation}

Thus, under $Z_2$ symmetry, the extrinsic curvature is fixed
\emph{algebraically} by the brane stress tensor. This is precisely why
the quadratic correction below descends to brane data while the projected
Weyl tensor does not.

\subsection{Gauss--Codazzi projection}
\label{sec:gauss}

The Gauss equation for the embedding is
\begin{equation}
\label{eq:gauss}
{}^{(4)}R_{\mu\nu\rho\sigma}
={}^{(5)}R_{ABCD}\,
e^A{}_\mu e^B{}_\nu e^C{}_\rho e^D{}_\sigma
+K_{\mu\rho}K_{\nu\sigma}-K_{\mu\sigma}K_{\nu\rho}.
\end{equation}
Decompose the five-dimensional Riemann tensor into Weyl and Ricci parts,
\begin{equation}
\label{eq:weyldecomp}
{}^{(5)}R_{ABCD}
={}^{(5)}C_{ABCD}
+\tfrac23\bigl(g_{A[C}{}^{(5)}R_{D]B}
-g_{B[C}{}^{(5)}R_{D]A}\bigr)
-\tfrac16\,{}^{(5)}R\,g_{A[C}g_{D]B},
\end{equation}
and use \eqref{eq:bulkEE}, which gives
${}^{(5)}R_{AB}=\tfrac23\Lfive g_{AB}$ and
${}^{(5)}R=\tfrac{10}{3}\Lfive$ away from the brane. Define the electric
projection of the bulk Weyl tensor,
\begin{equation}
\label{eq:Edef}
E_{\mu\nu}:={}^{(5)}C_{ABCD}\,
n^Ae^B{}_\mu\,n^Ce^D{}_\nu ,
\end{equation}
which is symmetric and traceless, $q^{\mu\nu}E_{\mu\nu}=0$, by the trace
identities of the Weyl tensor and $n\cdot n=1$.

Contracting \eqref{eq:gauss} twice and substituting
\eqref{eq:weyldecomp} yields the well-known intermediate identity
\begin{equation}
\label{eq:Gintermediate}
G_{\mu\nu}
=-\tfrac12\Lfive q_{\mu\nu}
+Q_{\mu\nu}(K)-E_{\mu\nu},
\end{equation}
where the complete extrinsic-curvature contribution is
\begin{equation}
\label{eq:QK}
Q_{\mu\nu}(K)
=KK_{\mu\nu}-K_{\mu\beta}K_\nu{}^\beta
-\tfrac12q_{\mu\nu}\bigl(K^2-K_{\alpha\beta}K^{\alpha\beta}\bigr).
\end{equation}

We now insert \eqref{eq:Kfixed} into \eqref{eq:QK}. Write
$K_{\mu\nu}=-\tfrac{\kfive^2}{2}\bigl(T_{\mu\nu}+\alpha q_{\mu\nu}\bigr)$
with $\alpha\equiv\tfrac13(\lambda-T)$, so that
$K=-\tfrac{\kfive^2}{2}(T+4\alpha)$. Then
\begin{align}
KK_{\mu\nu}
&=\frac{\kfive^4}{4}\bigl(T+4\alpha\bigr)
\bigl(T_{\mu\nu}+\alpha q_{\mu\nu}\bigr),
\label{eq:KK1}\\
K_{\mu\beta}K_\nu{}^\beta
&=\frac{\kfive^4}{4}
\bigl(T_{\mu\beta}T_\nu{}^\beta
+2\alpha T_{\mu\nu}+\alpha^2q_{\mu\nu}\bigr),
\label{eq:KK2}\\
K^2
&=\frac{\kfive^4}{4}\bigl(T+4\alpha\bigr)^2,
\qquad
K_{\alpha\beta}K^{\alpha\beta}
=\frac{\kfive^4}{4}
\bigl(T_{\alpha\beta}T^{\alpha\beta}
+2\alpha T+4\alpha^2\bigr).
\label{eq:KK3}
\end{align}
Substituting \eqref{eq:KK1}--\eqref{eq:KK3} directly into \eqref{eq:QK},
without yet using the value of $\alpha$, gives the unsimplified expression
\begin{equation}
\label{eq:Qexpanded}
\begin{aligned}
Q_{\mu\nu}
=\frac{\kfive^4}{4}\Bigl[
&\bigl(T+4\alpha\bigr)\bigl(T_{\mu\nu}+\alpha q_{\mu\nu}\bigr)
-\bigl(T_{\mu\beta}T_\nu{}^\beta
+2\alpha T_{\mu\nu}+\alpha^2q_{\mu\nu}\bigr)\\
&-\tfrac12 q_{\mu\nu}
\bigl(\bigl(T+4\alpha\bigr)^2
-T_{\alpha\beta}T^{\alpha\beta}
-2\alpha T-4\alpha^2\bigr)
\Bigr].
\end{aligned}
\end{equation}
Expanding and collecting on the independent structures
$T_{\mu\nu}$, $T_{\mu\beta}T_\nu{}^\beta$, and $q_{\mu\nu}$:
\begin{equation}
\label{eq:Qsimplified}
Q_{\mu\nu}
=\frac{\kfive^4}{4}\Bigl[
\bigl(T+2\alpha\bigr)T_{\mu\nu}
-T_{\mu\beta}T_\nu{}^\beta
+q_{\mu\nu}\Bigl(
\tfrac12 T_{\alpha\beta}T^{\alpha\beta}
-\tfrac12 T^2
-2\alpha T
-3\alpha^2
\Bigr)\Bigr].
\end{equation}
Indeed, the $q_{\mu\nu}$ coefficient collects
$\alpha(T+4\alpha)-\alpha^2
-\tfrac12\bigl(T^2+8\alpha T+16\alpha^2
-T_{\alpha\beta}T^{\alpha\beta}-2\alpha T-4\alpha^2\bigr)
=\tfrac12 T_{\alpha\beta}T^{\alpha\beta}-\tfrac12T^2-2\alpha T-3\alpha^2$,
while the $T_{\mu\nu}$ coefficient collects $(T+4\alpha)-2\alpha=T+2\alpha$.

Now restore $\alpha=\tfrac13(\lambda-T)$ in \eqref{eq:Qsimplified} and
collect by powers of $\lambda$. The $T_{\mu\nu}$ coefficient becomes
$T+\tfrac23(\lambda-T)=\tfrac13T+\tfrac23\lambda$, contributing
$\tfrac{\kfive^4}{4}\bigl(\tfrac13 T T_{\mu\nu}
+\tfrac23\lambda T_{\mu\nu}\bigr)
=\tfrac{\kfive^4}{12}TT_{\mu\nu}+\tfrac{\kfive^4\lambda}{6}T_{\mu\nu}$.
The $q_{\mu\nu}$ coefficient becomes
$\tfrac12 T_{\alpha\beta}T^{\alpha\beta}-\tfrac12T^2
-\tfrac23(\lambda-T)T-\tfrac13(\lambda-T)^2
=\tfrac12T_{\alpha\beta}T^{\alpha\beta}-\tfrac16T^2-\tfrac13\lambda^2$,
contributing
$\tfrac{\kfive^4}{8}q_{\mu\nu}T_{\alpha\beta}T^{\alpha\beta}
-\tfrac{\kfive^4}{24}q_{\mu\nu}T^2
-\tfrac{\kfive^4\lambda^2}{12}q_{\mu\nu}$. Together with the
$-\tfrac{\kfive^4}{4}T_{\mu\beta}T_\nu{}^\beta$ term this is exactly
\begin{equation}
\label{eq:Qsorted}
\boxed{\;
Q_{\mu\nu}
=-\frac{\kfive^4\lambda^2}{12}\,q_{\mu\nu}
+\frac{\kfive^4\lambda}{6}\,T_{\mu\nu}
+\kfive^4\,\Pi_{\mu\nu},
\;}
\end{equation}
with
\begin{equation}
\label{eq:Pi}
\boxed{\;
\Pi_{\mu\nu}
=-\tfrac14 T_{\mu\beta}T_\nu{}^\beta
+\tfrac1{12}\,T\,T_{\mu\nu}
+\tfrac18\,q_{\mu\nu}T_{\alpha\beta}T^{\alpha\beta}
-\tfrac1{24}\,q_{\mu\nu}T^2 .
\;}
\end{equation}
No cross terms are dropped and no $\lambda$-dependence is hidden in the
purely quadratic part: all $\lambda^2$ dependence enters through
$-\tfrac13\lambda^2$ in the $q_{\mu\nu}$ coefficient, and all $\lambda T$
dependence through $\tfrac23\lambda$ in the $T_{\mu\nu}$ coefficient.

Substituting \eqref{eq:Qsorted} into \eqref{eq:Gintermediate} gives the
Shiromizu--Maeda--Sasaki equation \cite{SMS}:
\begin{equation}
\label{eq:SMS}
\boxed{\;
G_{\mu\nu}[q]
=-\Lfour\,q_{\mu\nu}
+\kfour^2\,T_{\mu\nu}
+\kfive^4\,\Pi_{\mu\nu}
-E_{\mu\nu},
\;}
\end{equation}
with
\begin{equation}
\label{eq:constants}
\boxed{\;
\Lfour=\tfrac12\Bigl(\Lfive+\tfrac16\kfive^4\lambda^2\Bigr),
\qquad
\kfour^2=8\pi G_4=\tfrac16\kfive^4\lambda .
\;}
\end{equation}
Equivalently, writing the equation in the form
$G_{\mu\nu}+\Lfour q_{\mu\nu}=\kfour^2T_{\mu\nu}
+\mathcal C^{\RS}_{\mu\nu}$, the RS2 realization of the
source-induced metric response is
\begin{equation}
\label{eq:CRS}
\boxed{\;
\mathcal C^{\RS}_{\mu\nu}
=\kfive^4\,\Pi_{\mu\nu}-E_{\mu\nu}.
\;}
\end{equation}

Every coefficient in \eqref{eq:SMS}--\eqref{eq:CRS} follows from the
action \eqref{eq:action}, the junction condition \eqref{eq:israel}, and
the Gauss identity \eqref{eq:gauss}. Dimensions are consistent:
$[\kfive^4]=L^6$, $[T_{\mu\nu}T^{\mu\nu}]=L^{-8}$, so
$[\kfive^4\Pi]=L^{-2}=[E_{\mu\nu}]=[G_{\mu\nu}]$.

\begin{remark}[The superscript is not decorative]
The tensor \eqref{eq:CRS} is \emph{not} predicted by any abstract
source/readout considerations. It is the response obtained after choosing
five-dimensional Einstein gravity, the RS2 brane action, the $Z_2$
junction condition, and the stated source domain. Other source models
induce different response tensors.
\end{remark}

\subsection{Conservation, Bianchi, and nonclosure}
\label{sec:conservation}

With no exchange of matter energy into the bulk, the Codazzi equation
$D^\mu\bigl(K_{\mu\nu}-Kq_{\mu\nu}\bigr)
={}^{(5)}R_{AB}n^Ae^B{}_\nu=\tfrac23\Lfive\,n_Ae^A{}_\nu=0$
combined with \eqref{eq:Kfixed} gives
\begin{equation}
\label{eq:Tcons}
\nabla^\mu T_{\mu\nu}=0 .
\end{equation}
Applying the four-dimensional contracted Bianchi identity
$\nabla^\mu G_{\mu\nu}=0$ to \eqref{eq:SMS} then yields
\begin{equation}
\label{eq:bianchiE}
\boxed{\;
\nabla^\mu E_{\mu\nu}
=\kfive^4\,\nabla^\mu\Pi_{\mu\nu}.
\;}
\end{equation}

This is the exact nonclosure statement in local form. Equation
\eqref{eq:bianchiE} fixes the \emph{divergence} of $E_{\mu\nu}$, which is
four of its components' worth of information, but leaves its
transverse-traceless content --- precisely the free bulk gravitational
response --- undetermined by brane data. The projected system
\eqref{eq:SMS}--\eqref{eq:bianchiE} is covariantly consistent but is not
in general a closed four-dimensional initial-value system. This is the
tensorial counterpart of Proposition~\ref{prop:witness}: the wave packet
of Section~\ref{sec:witness} supplies two solutions with identical
instantaneous brane data and different later $E_{\mu\nu}$.

\section{Exact evaluation in the homogeneous sector}
\label{sec:frw}

Take an FRW brane and a perfect fluid,
\begin{equation}
\label{eq:frwansatz}
ds_4^2=-dt^2+a(t)^2\,d\Sigma_k^2,
\qquad
T_{\mu\nu}=\rho\,u_\mu u_\nu+p\,h_{\mu\nu},
\qquad
h_{\mu\nu}=q_{\mu\nu}+u_\mu u_\nu,
\end{equation}
with $u^\mu u_\mu=-1$. Then $T=-\rho+3p$ and
$T_{\alpha\beta}T^{\alpha\beta}=\rho^2+3p^2$.

\begin{lemma}[Perfect-fluid $\Pi_{\mu\nu}$]
\label{lem:PiFRW}
For \eqref{eq:frwansatz}, \eqref{eq:Pi} evaluates to
\begin{equation}
\label{eq:PiFRW}
\boxed{\;
\Pi_{\mu\nu}
=\frac{\rho}{12}\Bigl[\rho\,u_\mu u_\nu+(\rho+2p)\,h_{\mu\nu}\Bigr].
\;}
\end{equation}
\end{lemma}

\begin{proof}
Direct substitution. In the rest frame, $T_{00}=\rho$,
$T_{ij}=p\,q_{ij}$, and
\begin{equation}
\Pi_{00}
=\frac{\rho^2}{4}
+\frac{(-\rho+3p)\rho}{12}
-\frac{\rho^2+3p^2}{8}
+\frac{(-\rho+3p)^2}{24}
=\frac{\rho^2}{12},
\end{equation}
where the first term carries a $+$ sign because
$-\tfrac14T_{0\beta}T_0{}^\beta
=-\tfrac14 T_{00}T_0{}^0=-\tfrac14\rho\cdot(-\rho)=+\tfrac14\rho^2$,
and the third and fourth carry the signs induced by $q_{00}=-1$.
Assembling the terms:
$\tfrac{\rho^2}{4}
-\tfrac{\rho^2}{12}+\tfrac{p\rho}{4}
-\tfrac{\rho^2}{8}-\tfrac{3p^2}{8}
+\tfrac{\rho^2}{24}-\tfrac{p\rho}{4}+\tfrac{3p^2}{8}
=\rho^2\bigl(\tfrac14-\tfrac1{12}-\tfrac18+\tfrac1{24}\bigr)
=\tfrac{\rho^2}{12}$.
Similarly $\Pi_{ij}=\tfrac{1}{12}\rho(\rho+2p)\,q_{ij}$. These are the
$u_\mu u_\nu$ and $h_{\mu\nu}$ components of \eqref{eq:PiFRW}.
\end{proof}

Homogeneity, isotropy, tracelessness, and \eqref{eq:bianchiE} (whose
right-hand side vanishes in this sector because $\Pi_{\mu\nu}$ is a
conserved perfect-fluid combination) force
\begin{equation}
\label{eq:EFRW}
E_{\mu\nu}=-\mathcal U(t)\Bigl(u_\mu u_\nu+\tfrac13h_{\mu\nu}\Bigr),
\qquad
\dot{\mathcal U}+4H\mathcal U=0
\quad\Longrightarrow\quad
\mathcal U(t)=\frac{3\mu}{a^4(t)} ,
\end{equation}
where $\mu$ is a constant fixed by the bulk geometry --- in an
AdS--Schwarzschild realization, proportional to the bulk black-hole mass
parameter \cite{Mukohyama1,BDEL}. It is \emph{not} determined by the
brane matter state.

The response is now fully evaluated:
\begin{equation}
\label{eq:CFRW}
\boxed{\;
\mathcal C^{\RS}_{00}
=\frac{\kfive^4}{12}\rho^2+\frac{3\mu}{a^4},
\qquad
\mathcal C^{\RS}_{ij}
=\Bigl[\frac{\kfive^4}{12}\rho(\rho+2p)+\frac{\mu}{a^4}\Bigr]q_{ij}.
\;}
\end{equation}

\begin{proposition}[Exact brane Friedmann equations]
\label{prop:friedmann}
For the model of Section~\ref{sec:model} with an FRW brane, conserved
brane matter, and no bulk matter flux,
\begin{equation}
\label{eq:friedmann}
\boxed{\;
H^2+\frac{k}{a^2}
=\frac{\Lfour}{3}
+\frac{\kfour^2}{3}\rho
+\frac{\kfive^4}{36}\rho^2
+\frac{\mu}{a^4}.
\;}
\end{equation}
Using $\kfour^2=\kfive^4\lambda/6$ this is equivalently
\begin{equation}
\label{eq:friedmann2}
H^2+\frac{k}{a^2}
=\frac{\Lfour}{3}
+\frac{\kfour^2}{3}\rho\Bigl(1+\frac{\rho}{2\lambda}\Bigr)
+\frac{\mu}{a^4},
\end{equation}
and, with $\dot\rho+3H(\rho+p)=0$,
\begin{equation}
\label{eq:accel}
\dot H-\frac{k}{a^2}
=-\frac{\kfour^2}{2}(\rho+p)\Bigl(1+\frac{\rho}{\lambda}\Bigr)
-\frac{2\mu}{a^4}.
\end{equation}
\end{proposition}

\begin{proof}
Take the $00$ component of \eqref{eq:SMS} with
\eqref{eq:PiFRW}--\eqref{eq:EFRW}. The Einstein tensor gives
$G_{00}=3(H^2+k/a^2)$; the correction contributes
$\kfive^4\Pi_{00}-E_{00}
=\tfrac{\kfive^4}{12}\rho^2+3\mu/a^4$. Dividing by $3$ gives
\eqref{eq:friedmann}. Substituting
$\kfive^4=6\kfour^2/\lambda$ gives
$\kfive^4\rho^2/36=\kfour^2\rho^2/(6\lambda)$, i.e.\
the $\rho^2/(2\lambda)$ term in \eqref{eq:friedmann2}. Differentiating
\eqref{eq:friedmann2} and eliminating $\dot\rho$ with conservation gives
\eqref{eq:accel}.
\end{proof}

The signs and scalings are fixed and physically transparent. Positive
brane tension gives a \emph{positive} $\rho^2$ correction, enhancing the
expansion rate at high density. The Weyl contribution scales as
radiation, but its magnitude is set by a bulk-state parameter, not by
brane matter. For $\rho/\lambda\ll1$ and $\mu/a^4\ll H^2$, standard
four-dimensional cosmology is recovered.

The homogeneous sector already exhibits the exact division between what
the projection fixes and what it leaves open:
\begin{itemize}[leftmargin=*]
\item the form and coefficient of the $\rho^2$ term are fixed once
$\kfive$ and $\lambda$ are specified;
\item the exponent $a^{-4}$ is fixed by symmetry and conservation alone;
\item the constant $\mu$ is bulk initial/boundary data, undetermined by
the brane matter state. It is not a prediction of the model without a
bulk-state condition.
\end{itemize}

\section{Exact tree-level bulk elimination in the linear weak-field sector}
\label{sec:linear}

The homogeneous sector evaluates the nonlinear local response but hides
the spectral structure of the eliminated bulk gravitons. That structure
is explicit at linear order \cite{RS2,GarrigaTanaka,GKR}.

\subsection{Bulk perturbation equation, brane bending, and the boundary condition}
\label{sec:bending}

In Randall--Sundrum gauge, a TT perturbation of the background
\eqref{eq:background} obeys
\begin{equation}
\label{eq:bulkwave}
\Bigl[
e^{2|y|/\ell}\,\Box_4
+\partial_y^2-\frac{4}{\ell^2}
+\frac{4}{\ell}\,\delta(y)
\Bigr]h^{\TT}_{\mu\nu}
=-2\kfive^2\,\Sigma^{\TT}_{\mu\nu}\,\delta(y).
\end{equation}
The source $\Sigma_{\mu\nu}$ contains the brane stress tensor
\emph{and} the brane-bending contribution required by the junction
condition. In the Garriga--Tanaka convention \cite{GarrigaTanaka},
\begin{equation}
\label{eq:bending}
\boxed{\;
\Sigma_{\mu\nu}
=T_{\mu\nu}-\tfrac13\eta_{\mu\nu}T
+\frac{2}{\kfive^2}\,\partial_\mu\partial_\nu\hat\xi^5,
\qquad
\Box_4\hat\xi^5=\frac{\kfive^2}{6}\,T .
\;}
\end{equation}
The bending equation makes $\Sigma^\mu{}_\mu=0$: indeed
$\eta^{\mu\nu}\Sigma_{\mu\nu}
=T-\tfrac43T+\tfrac{2}{\kfive^2}\Box_4\hat\xi^5
=-\tfrac13T+\tfrac13T=0$.
Together with $\partial^\mu T_{\mu\nu}=0$ it also gives
$\partial^\mu\Sigma_{\mu\nu}=0$, as RS gauge requires.

Brane bending is not an optional refinement. It is the statement that
the brane, being a dynamical surface, is displaced from $y=0$ in RS gauge
by the presence of matter; restoring the Gaussian-normal frame adapted to
the physical brane produces, up to a tangential gauge transformation,
\begin{equation}
\label{eq:gaussnormal}
\bar h_{\mu\nu}
=h_{\mu\nu}
+\ell\,\partial_\mu\partial_\nu\hat\xi^5
+\frac{2}{\ell}\,\eta_{\mu\nu}\hat\xi^5
-\partial_{(\mu}\hat\xi_{\nu)} .
\end{equation}
Decomposing $h_{\mu\nu}=h^{(m)}_{\mu\nu}+h^{(\xi)}_{\mu\nu}$ into the
part sourced by $T_{\mu\nu}-\tfrac13\eta_{\mu\nu}T$ and the part sourced
by the bending term, and choosing the tangential gauge to cancel the pure
gradient pieces, the brane value reduces to
\begin{equation}
\label{eq:hbar}
\boxed{\;
\bar h_{\mu\nu}
=h^{(m)}_{\mu\nu}
+\frac{2}{\ell}\,\eta_{\mu\nu}\,\hat\xi^5 .
\;}
\end{equation}
This single term converts the apparent $\tfrac13$ trace coefficient of
the RS-gauge source into the four-dimensional massless $\tfrac12$
coefficient in the zero-mode response --- as Section~\ref{sec:projectors}
makes explicit. Omitting brane bending, or using the wrong trace
coefficient, produces a graviton that is not the four-dimensional
graviton and gives the wrong light-bending prediction.

The retarded Green function imposes causal response; regularity, i.e.\
absence of incoming radiation from the AdS horizon, selects the forced
solution. If arbitrary incoming bulk radiation is admitted, a homogeneous
solution must be added; \emph{that homogeneous term is exactly the
initial-state contribution associated with the same-readout fiber of
Section~\ref{sec:witness}}. The two constructions are the same physics
seen twice.

\subsection{Brane-to-brane propagator: zero mode and KK continuum}
\label{sec:propagator}

The brane-to-brane retarded Green function of \eqref{eq:bulkwave} has the
spectral resolution
\begin{equation}
\label{eq:GR}
\boxed{\;
\begin{aligned}
G_R(p;0,0)
&=G_R^{(0)}(p)+G_R^{(\KK)}(p),\\
G_R^{(0)}(p)
&=-\frac{1}{\ell}\,
\frac{1}{-(\omega+i0^+)^2+\mathbf k^2},\\
G_R^{(\KK)}(p)
&=-\int_0^\infty dm\,
\frac{u_m(0)^2}{m^2-(\omega+i0^+)^2+\mathbf k^2} .
\end{aligned}
\;}
\end{equation}
The normalized continuum modes of the transverse operator are Bessel
combinations, and their brane value gives the spectral weight
\begin{equation}
\label{eq:um}
\boxed{\;
u_m(0)^2
=\frac{2}{\pi^2 m\ell\bigl[J_1^2(m\ell)+Y_1^2(m\ell)\bigr]} .
\;}
\end{equation}
Using the small-argument behavior
$J_1(x)\simeq x/2$ and $Y_1(x)\simeq-2/(\pi x)$, one has
$J_1^2+Y_1^2\simeq 4/(\pi^2x^2)$, hence
\begin{equation}
\label{eq:umsmall}
u_m(0)^2
=\frac{m\ell}{2}
+O\bigl((m\ell)^3\log(m\ell)\bigr),
\qquad m\ell\ll1 .
\end{equation}

Two structural facts follow. The first term in \eqref{eq:GR} is the
normalizable four-dimensional graviton zero mode: an isolated massless
pole. The second is a \emph{continuum} of physical KK modes with
nonnegative spectral density $u_m(0)^2\ge0$, producing a branch cut
beginning at threshold $m=0$.

\begin{remark}[What the branch cut does and does not exclude]
\label{rem:cut}
A finite-dimensional classical linear system, or a finite set of free
fields with a finite-order quadratic action evaluated at tree level, has
only a finite rational pole structure and cannot reproduce
\eqref{eq:GR}--\eqref{eq:um} exactly. It does \emph{not} follow that
every four-dimensional theory is excluded. An interacting four-dimensional
field theory, a nonlocal theory, an infinite tower, a continuum of
fields, or a conformal sector reproduces a branch cut without difficulty
--- multiparticle thresholds generically produce cuts. See
Sections~\ref{sec:fourdim} and \ref{sec:failure}, item 8.
\end{remark}

\subsection{Tensor projectors and the response operator}
\label{sec:projectors}

At linear order the bulk gravitational action is quadratic. Integration
over bulk modes is therefore Gaussian, and the resulting Schur complement
is \emph{exact at tree level}. We first retain the tensor structure that
distinguishes the massless graviton from the massive continuum.

For four-momentum $p$, define the transverse tensor
\begin{equation}
\label{eq:theta}
\theta_{\mu\nu}
=\eta_{\mu\nu}-\frac{p_\mu p_\nu}{p^2+i0^+},
\end{equation}
and the Barnes--Rivers spin projectors on symmetric tensors
\cite{PinheiroPires},
\begin{equation}
\label{eq:BR}
\begin{aligned}
P^{(2)}_{\mu\nu\rho\sigma}
&=\tfrac12\bigl(\theta_{\mu\rho}\theta_{\nu\sigma}
+\theta_{\mu\sigma}\theta_{\nu\rho}\bigr)
-\tfrac13\theta_{\mu\nu}\theta_{\rho\sigma},\\
P^{(0)}_{\mu\nu\rho\sigma}
&=\tfrac13\theta_{\mu\nu}\theta_{\rho\sigma}.
\end{aligned}
\end{equation}
(The $\tfrac13$ is $1/(D-1)$ with $D=4$, appropriate to the
three-dimensional transverse subspace.)

\begin{lemma}[Trace coefficients]
\label{lem:traces}
For a conserved source, $p^\mu T_{\mu\nu}=0$, so
$\theta_{\mu\nu}T^{\mu\nu}=T$ and
\begin{equation}
\label{eq:projT}
\bigl(P^{(2)}T\bigr)_{\mu\nu}=T_{\mu\nu}-\tfrac13\theta_{\mu\nu}T,
\qquad
\bigl(P^{(0)}T\bigr)_{\mu\nu}=\tfrac13\theta_{\mu\nu}T .
\end{equation}
Consequently
\begin{equation}
\label{eq:massless}
\Bigl(P^{(2)}-\tfrac12P^{(0)}\Bigr)T
=T_{\mu\nu}-\tfrac12\theta_{\mu\nu}T ,
\end{equation}
which is the massless four-dimensional graviton structure, while a
massive spin-two exchange gives $P^{(2)}T=T_{\mu\nu}
-\tfrac13\theta_{\mu\nu}T$, the van Dam--Veltman--Zakharov structure.
\end{lemma}

\begin{proof}
Contract \eqref{eq:BR} with $T^{\rho\sigma}$ using
$\theta_{\mu\rho}T^{\rho\sigma}\theta_{\sigma\nu}=T_{\mu\nu}$ for a
conserved source, and $\theta_{\rho\sigma}T^{\rho\sigma}=T$. Then
$P^{(2)}T=T_{\mu\nu}-\tfrac13\theta_{\mu\nu}T$ and
$P^{(0)}T=\tfrac13\theta_{\mu\nu}T$. Subtracting half the latter,
$T_{\mu\nu}-\tfrac13\theta_{\mu\nu}T-\tfrac16\theta_{\mu\nu}T
=T_{\mu\nu}-\tfrac12\theta_{\mu\nu}T$.
\end{proof}

The brane-bending field of Section~\ref{sec:bending} is precisely what
converts the scalar coefficient of the massless part from $\tfrac13$ to
$\tfrac12$: the term $\tfrac{2}{\ell}\eta_{\mu\nu}\hat\xi^5$ in
\eqref{eq:hbar} supplies a trace piece proportional to
$\Box_4^{-1}T$, i.e.\ $-\tfrac16\theta_{\mu\nu}T$ times the zero-mode
propagator. Consequently the complete gauge-invariant response to a
conserved brane source, modulo longitudinal pure-gauge terms, is
\begin{equation}
\label{eq:response}
\boxed{\;
\bar h_{\mu\nu}(p)
=-\mathbb R_{R\,\mu\nu}{}^{\rho\sigma}(p)\,T_{\rho\sigma}(p),
\;}
\end{equation}
with
\begin{equation}
\label{eq:RR}
\boxed{\;
\mathbb R_R(p)
=2\kfive^2\Bigl[
G_R^{(0)}(p)\Bigl(P^{(2)}-\tfrac12P^{(0)}\Bigr)
+G_R^{(\KK)}(p)\,P^{(2)}
\Bigr].
\;}
\end{equation}
Equation \eqref{eq:RR} fixes both the normalization and the relative
tensor coefficients. It is equivalent to the RS-gauge source equation
\eqref{eq:bulkwave} together with the bending equation
\eqref{eq:bending} and the transformation \eqref{eq:hbar}. It displays
directly why the zero mode carries the four-dimensional massless trace
coefficient while the KK continuum carries the massive spin-two
coefficient --- the physical content of the fact that RS2 has no vDVZ
discontinuity in the observable weak-field limit.

\subsection{Causal influence functional}
\label{sec:ctp}

The same result follows from a causal influence functional, which is the
formulation that makes the initial-state dependence explicit
\cite{FeynmanVernon}. Causality requires the closed-time-path fields to
remain doubled. Set
\begin{equation}
\label{eq:keldysh}
h_c=\tfrac12(h_++h_-),
\qquad
h_\Delta=h_+-h_- .
\end{equation}
On the physical conserved-source subspace let
$\mathbb K_R=\mathbb R_R^{-1}$. For equal physical sources on the two
contour branches, the quadratic closed-time-path functional is
\begin{equation}
\label{eq:CTP}
\boxed{\;
\Gamma^{(2)}_{\mathrm{CTP}}[h_c,h_\Delta;T]
=\tfrac12\bigl\langle h_\Delta,\mathbb K_R h_c\bigr\rangle
+\tfrac12\bigl\langle h_\Delta,T\bigr\rangle
+\tfrac{i}{2}\bigl\langle h_\Delta,\mathbb N h_\Delta\bigr\rangle ,
\;}
\end{equation}
where the pairing includes momentum integration and contraction on the
conserved tensor subspace. The factor $\tfrac12$ in the source term is
the standard linear metric coupling
$\delta S_m=\tfrac12\int T^{\mu\nu}\delta h_{\mu\nu}$. The symmetric
kernel $\mathbb N$ is fixed by the chosen bulk density matrix; it
controls fluctuations and does not enter the mean classical field.
Varying \eqref{eq:CTP} with respect to $h_\Delta$ and then taking the
physical limit $h_\Delta=0$ gives
\begin{equation}
\label{eq:ctpEOM}
\tfrac12\,\mathbb K_R h_c+\tfrac12\,T=0
\qquad\Longrightarrow\qquad
h_c=-\mathbb R_R\,T ,
\end{equation}
reproducing \eqref{eq:response}. Restricted to the TT source in RS gauge
this reduces to
\begin{equation}
\label{eq:TTresponse}
h_c^{\TT}=-2\kfive^2\,\mathcal G_R\,\Sigma^{\TT},
\end{equation}
where $\mathcal G_R$ denotes the retarded convolution operator whose
momentum-space kernel is the brane-to-brane Green function
$G_R(p;0,0)=G_R^{(0)}(p)+G_R^{(\KK)}(p)$ of \eqref{eq:GR}, acting as
\begin{equation}
\label{eq:GRop}
\bigl(\mathcal G_R\,\Sigma^{\TT}\bigr)(x)
=\int d^4x'\;\mathcal G_R(x-x')\,\Sigma^{\TT}(x'),
\qquad
\mathcal G_R(x-x')
=\int\frac{d^4p}{(2\pi)^4}\;e^{ip\cdot(x-x')}\,G_R(p;0,0).
\end{equation}
This recovers the normalization used in \eqref{eq:bulkwave}. Thus the
standard $\tfrac12 h_{\mu\nu}T^{\mu\nu}$ matter coupling \emph{reproduces}
the response normalization rather than fixing it by convention.

The retarded prescription fixes causality. In the positive-tension model
the KK spectral density \eqref{eq:um} is positive, so the linear
continuum introduces no negative-norm graviton. The negative-tension
branch instead gives the wrong sign of the four-dimensional Newton
constant and is not used.

\begin{remark}[Status of the elimination]
\label{rem:treelevel}
This is a \emph{tree-level, linear-response} calculation, and must not be
confused with the exact nonlinear equation \eqref{eq:SMS}. Gauge fixing
isolates the physical TT operator; brane bending restores the response
required by the conserved four-dimensional source and its Ward identity.
At tree level, ghost fields carry no external classical source and
contribute no exchange term between conserved brane stresses; ghost and
graviton determinants enter only beyond this approximation. Ultraviolet
local terms at coincident points renormalize the brane tension, the
Newton coupling, and admissible higher-curvature operators; the
nonanalytic continuum term and its $1/r^3$ residue at $r>0$ are
insensitive to those local subtractions. Covariant conservation of the
resulting response agrees with \eqref{eq:bianchiE}. The bulk initial
state enters the full closed-time-contour functional through its density
matrix: under the no-incoming/regular choice the forced retarded kernel
is fixed; changing that choice changes both the homogeneous response and
the fluctuation kernel $\mathbb N$.
\end{remark}

\subsection{Evaluated long-distance response}
\label{sec:longdistance}

For a static source and $r\gg\ell$, the three-dimensional massive Green
function is $-e^{-mr}/(4\pi r)$. The small-$m$ spectral weight
\eqref{eq:umsmall} therefore \emph{evaluates} the KK contribution rather
than leaving it as a formal kernel:
\begin{equation}
\label{eq:GKKr}
G^{(\KK)}(r)
=-\frac{1}{4\pi r}\int_0^\infty dm\,u_m(0)^2\,e^{-mr}
=-\frac{1}{4\pi r}\cdot\frac{\ell}{2}\int_0^\infty dm\,m\,e^{-mr}
+\cdots
=-\frac{\ell}{8\pi r^3}
+O\!\Bigl(\frac{\ell^3}{r^5}\log\frac r\ell\Bigr),
\end{equation}
using $\int_0^\infty m e^{-mr}dm=r^{-2}$. Adding the zero mode
$G^{(0)}(r)=-1/(4\pi\ell r)$ gives the brane scalar Green function
\begin{equation}
\label{eq:Gtotal}
G(\mathbf x,0;\mathbf x',0)
=-\frac{1}{4\pi\ell r}
\Bigl[1+\frac{\ell^2}{2r^2}
+O\!\Bigl(\frac{\ell^4}{r^4}\log\frac r\ell\Bigr)\Bigr],
\qquad r=|\mathbf x-\mathbf x'| .
\end{equation}

For a nonrelativistic point source, $T_{00}=M\delta^3(\mathbf x)$ and
$T=-M\delta^3(\mathbf x)$, so the bending equation \eqref{eq:bending}
gives
\begin{equation}
\label{eq:xi5}
\nabla^2\hat\xi^5=-\frac{\kfive^2M}{6}\,\delta^3(\mathbf x),
\qquad
\hat\xi^5(r)=\frac{\kfive^2M}{24\pi r} .
\end{equation}
Substituting \eqref{eq:Gtotal} into
$h^{(m)}=-2\kfive^2\,G*\bigl(T_{\mu\nu}-\tfrac13\eta_{\mu\nu}T\bigr)$,
adding $2\eta_{\mu\nu}\hat\xi^5/\ell$ as in \eqref{eq:hbar}, and using
$G_4=\kfive^2/(8\pi\ell)$ (Section~\ref{sec:normalization}), one obtains,
outside a nonrelativistic spherical source of mass $M$,
\begin{equation}
\label{eq:h00}
\boxed{\;
\bar h_{00}
=\frac{2G_4M}{r}
\Bigl[1+\frac{2\ell^2}{3r^2}
+O\!\Bigl(\frac{\ell^4}{r^4}\log\frac r\ell\Bigr)\Bigr],
\;}
\end{equation}
\begin{equation}
\label{eq:hij}
\boxed{\;
\bar h_{ij}
=\frac{2G_4M}{r}
\Bigl[1+\frac{\ell^2}{3r^2}
+O\!\Bigl(\frac{\ell^4}{r^4}\log\frac r\ell\Bigr)\Bigr]\delta_{ij}.
\;}
\end{equation}
Equivalently, in asymptotically Minkowski Newtonian gauge,
\begin{equation}
\label{eq:newtoniangauge}
ds^2=-(1+2\Phi)\,dt^2+(1-2\Psi)\,d\mathbf x^2,
\end{equation}
the gauge-invariant static potentials are
\begin{equation}
\label{eq:potentials}
\boxed{\;
\Phi(r)=-\frac{G_4M}{r}
\Bigl(1+\frac{2\ell^2}{3r^2}+\cdots\Bigr),
\qquad
\Psi(r)=-\frac{G_4M}{r}
\Bigl(1+\frac{\ell^2}{3r^2}+\cdots\Bigr).
\;}
\end{equation}
This reproduces the Garriga--Tanaka result \cite{GarrigaTanaka}.

Three observations. First, the correction is \emph{not} a mass
renormalization: it has a different radial scaling and cannot be absorbed
into $G_4M$. Second, the coefficients in $\Phi$ and $\Psi$ are unequal,
which correlates nonrelativistic acceleration with light deflection: the
former probes $\Phi$, the latter $\Phi+\Psi$. Third, in momentum space
the same physics is the KK branch cut of \eqref{eq:GR}, not an extra
isolated pole.

The result is controlled by
\begin{equation}
\label{eq:controlconds}
\frac{G_4M}{r}\ll1,
\qquad
\frac{\ell}{r}\ll1 ,
\end{equation}
and assumes the minimal two-derivative bulk action and the regular
retarded bulk state. \emph{It is not an expansion valid at
$r\lesssim\ell$}, where the full spectral integral \eqref{eq:um} must be
used.

\subsection{One-sided versus two-sided normalization}
\label{sec:normalization}

Two conventions occur in the literature and must not be mixed. In the
two-sided covering-space convention used throughout this paper,
\begin{equation}
\label{eq:twosided}
y\in(-\infty,\infty),
\qquad
a(y)=e^{-|y|/\ell},
\end{equation}
and the zero-mode reduction gives
\begin{equation}
\label{eq:kappa4}
\frac{1}{\kfour^2}
=\frac{1}{\kfive^2}\int_{-\infty}^{\infty}dy\;a(y)^2
=\frac{1}{\kfive^2}\int_{-\infty}^{\infty}dy\;e^{-2|y|/\ell}
=\frac{\ell}{\kfive^2},
\end{equation}
so that
\begin{equation}
\label{eq:G4}
\boxed{\;
\kfour^2=\frac{\kfive^2}{\ell},
\qquad
G_4=\frac{G_5}{\ell}.
\;}
\end{equation}
If instead one works on a single AdS patch $y\ge0$ with the same
undoubled bulk normalization, the integral is $\ell/2$ and
\begin{equation}
\label{eq:onesided}
\kfour^2=\frac{2\kappa_{5,\mathrm{one}}^2}{\ell},
\qquad
G_4=\frac{2G_{5,\mathrm{one}}}{\ell}.
\end{equation}
Equivalently, the $Z_2$ orbifold fundamental domain may be assigned a
doubled bulk action, reproducing the two-sided convention. The weak-field
coefficients $2\ell^2/3$ and $\ell^2/3$, expressed in terms of the
\emph{measured} $G_4$, are unchanged provided the gravitational coupling
and the number of AdS patches are converted consistently.

The same distinction appears in the four-dimensional dual description
(Section~\ref{sec:fourdim}). One AdS patch corresponds to one cutoff CFT
contribution; a two-sided geometry has two identical bulk regions and
hence a doubled CFT polarization contribution. Duff and Liu \cite{DuffLiu}
present the one-patch normalization explicitly; doubling both the bulk
contribution and its CFT counterpart gives the two-sided convention used
here. Their equality of the nonanalytic $1/r^3$ residue therefore remains
valid, but only after this conversion. No equation in this paper combines
a one-sided Newton-constant relation with a two-sided junction condition.

\section{The parameter-eliminated cross-regime prediction}
\label{sec:crossregime}

Under the RS tuning
\begin{equation}
\label{eq:tuning}
\Lfive=-\frac{6}{\ell^2},
\qquad
\lambda=\frac{6}{\kfive^2\ell},
\qquad
\kfour^2=\frac{\kfive^2}{\ell},
\end{equation}
one also has $\lambda=6/(\kfour^2\ell^2)$, and \eqref{eq:constants} gives
$\Lfour=\tfrac12\bigl(-6/\ell^2+\tfrac16\kfive^4\lambda^2\bigr)=0$
exactly, since $\kfive^4\lambda^2=\kfive^4\cdot36/(\kfive^4\ell^2)
=36/\ell^2$.

\begin{remark}[Status of $\Lfour$ below]
The symbol $\Lfour$ retained in \eqref{eq:phenomB} is a
\emph{phenomenological} background term that may be fitted or supplied by
a much smaller late-time brane energy component. It is not an independent
parameter of the exactly tuned flat-brane theorem. Applying the
coefficient relation on a weakly curved cosmological background requires
$H\ell\ll1$ and negligible detuning corrections in the compared regime.
See Section~\ref{sec:failure}, item 9.
\end{remark}

Write the two independently measurable departures from four-dimensional
GR, \emph{without} using RS2 in the fits, as
\begin{equation}
\label{eq:phenomB}
H^2+\frac{k}{a^2}
=\frac{\Lfour}{3}+\frac{\kfour^2}{3}\rho
+B\rho^2+\frac{\mu}{a^4},
\end{equation}
\begin{equation}
\label{eq:phenomA}
\Phi(r)=-\frac{G_4M}{r}\Bigl(1+\frac{A_\Phi}{r^2}+\cdots\Bigr),
\qquad
\Psi(r)=-\frac{G_4M}{r}\Bigl(1+\frac{A_\Psi}{r^2}+\cdots\Bigr).
\end{equation}

\begin{theorem}[Cross-regime coefficient relation]
\label{thm:crossregime}
In minimal, exactly tuned, positive-tension RS2 with the regular retarded
bulk state, the source calculation gives
\begin{equation}
\label{eq:sourcecoeffs}
B=\frac{\kfive^4}{36}=\frac{\kfour^4\ell^2}{36},
\qquad
A_\Phi=\frac{2\ell^2}{3},
\qquad
A_\Psi=\frac{\ell^2}{3},
\end{equation}
and eliminating $\ell$ yields
\begin{equation}
\label{eq:crossregime}
\boxed{\;
A_\Psi=\tfrac12 A_\Phi,
\qquad
B=\frac{\kfour^4}{24}A_\Phi
=\frac{\kfour^4}{12}A_\Psi .
\;}
\end{equation}
\end{theorem}

\begin{proof}
$B=\kfive^4/36$ is read off from \eqref{eq:friedmann}. Using
$\kfive^2=\kfour^2\ell$ from \eqref{eq:G4} gives
$\kfive^4=\kfour^4\ell^2$, hence $B=\kfour^4\ell^2/36$. The values
$A_\Phi=2\ell^2/3$, $A_\Psi=\ell^2/3$ are read off from
\eqref{eq:potentials}. Eliminating, $\ell^2=3A_\Phi/2=3A_\Psi$, so
$A_\Psi=A_\Phi/2$, and
$B=\tfrac{\kfour^4}{36}\cdot\tfrac{3A_\Phi}{2}
=\tfrac{\kfour^4}{24}A_\Phi
=\tfrac{\kfour^4}{12}A_\Psi$.
\end{proof}

\paragraph{Dimensional consistency.}
$[\kfour^2]=L^2$, $[A_\Phi]=L^2$, hence $[B]=L^4\cdot L^2=L^6$; and with
$[\rho]=L^{-4}$, $[B\rho^2]=L^{6-8}=L^{-2}=[H^2]$, as required.

\paragraph{Scope, and what it does and does not do.}
The relation \eqref{eq:crossregime} contains no unknown brane parameter
once $G_4$ and the phenomenological coefficients are measured. It is
independent of the dark-radiation constant $\mu$, whose $a^{-4}$ scaling
is separable from $\rho^2$ when the matter equation of state and energy
density are independently constrained. Its content is exact within its
\emph{intersecting} regimes:
\begin{itemize}[leftmargin=*]
\item $B$ comes from the exact homogeneous thin-brane equation
\eqref{eq:friedmann};
\item $A_\Phi,A_\Psi$ are the leading terms of the controlled
$r\gg\ell$ weak-field expansion \eqref{eq:potentials};
\item the same $\ell$ and the same tuned source action govern both
sectors;
\item no fit of $\ell$ remains after the coefficients are compared.
\end{itemize}

A violation of either equality falsifies minimal classical RS2 under the
stated assumptions. A statistically significant \emph{nonzero} common
signal satisfying both equalities is a parameter-saving, out-of-sample
success for RS2 and favours it over pure four-dimensional GR in model
comparison. It does \emph{not} falsify a four-dimensional theory that
independently parameterizes all three coefficients: such a theory fits
either agreement or violation. Only a restricted alternative predicting
an \emph{incompatible} relation is excluded. Note also the all-zero limit:
$A_\Phi=A_\Psi=B=0$ satisfies both equalities trivially and merely bounds
$\ell$ from above. And even nonzero agreement does not establish that the
five-dimensional variables are ontologically fundamental, as
Section~\ref{sec:opeq} proves.

\section{Minimal dynamical completion and memory}
\label{sec:completion}

\subsection{Linear completion: a quotient by the kernel of the brane response}

Restrict first to the linearized, fixed-background problem with fixed
brane forcing, for which the gauge-reduced bulk perturbation data form a
vector space. Let $z_0$ denote the bulk gravitational initial data
omitted by the readout \eqref{eq:readout}, and define the brane-response
operator
\begin{equation}
\label{eq:Osigma}
\mathcal O_\Sigma z_0
=\text{the induced brane response generated by }z_0,
\end{equation}
with brane initial data and applied brane sources held fixed. Put
$N_\Sigma=\ker\mathcal O_\Sigma$.

\begin{definition}[Minimal linear dynamical completion]
\label{def:Edyn}
\begin{equation}
\label{eq:Edyn}
\boxed{\;
E^{\RS}_{\mathrm{dyn}}
=\frac{\{\text{admissible gauge-reduced bulk initial data}\}}
{N_\Sigma} .
\;}
\end{equation}
\end{definition}

This quotient discards exactly the bulk directions that never influence
the brane, and retains no more. In the linear spectral representation it
retains precisely the coupled KK amplitudes together with the allowed
incoming homogeneous modes. Under a prescribed regular no-incoming state,
the homogeneous amplitudes are fixed and the entire remaining response is
the retarded convolution encoded by $G_R^{(\KK)}$.

Schematically, in each relevant projected tensor sector, the exact linear
brane equation takes the causal, generally nonlocal form
\begin{equation}
\label{eq:memoryeq}
\mathcal D_0 h_\Sigma(x)
+\int d^4x'\;\mathcal K_R(x,x')\,h_\Sigma(x')
=T_\Sigma(x)+\eta_{z_0}(x),
\end{equation}
with
\begin{equation}
\label{eq:KR}
\mathcal K_R
=\bigl(G_R^{(0)}+G_R^{(\KK)}\bigr)^{-1}
-\bigl(G_R^{(0)}\bigr)^{-1}.
\end{equation}
Tensor indices and normalizations are suppressed throughout
\eqref{eq:memoryeq}--\eqref{eq:KR}: the symbols $\mathcal D_0$,
$T_\Sigma$, and $\mathcal K_R$ denote the appropriately normalized
operators and sources within the projected tensor sector under
consideration. In equation \eqref{eq:KR},
$G_R^{(0)}+G_R^{(\KK)}$ is specifically the scalar coefficient
multiplying $P^{(2)}$ in the full response operator \eqref{eq:RR}; thus
$\mathcal K_R$ is the spin-two-sector correction relative to the
zero-mode operator. The $P^{(0)}$ contribution occurs only through the
massless combination $G_R^{(0)}\bigl(P^{(2)}-\tfrac12P^{(0)}\bigr)$ and
must be retained through the full projector-valued response
\eqref{eq:RR}. The combination $P^{(2)}-\tfrac12P^{(0)}$ is the complete
massless-graviton exchange structure, not a scalar sector. Equation
\eqref{eq:memoryeq} is accordingly a schematic sector-by-sector statement,
not a fully explicit tensor identity; the explicit tensor content is
\eqref{eq:response}--\eqref{eq:RR}. The term $\eta_{z_0}$ is the
homogeneous response of the unresolved bulk initial data. Two points
deserve emphasis. First, $\mathcal K_R$ is
\emph{computed} by the KK spectral integral \eqref{eq:um}; it is not an
arbitrary phenomenological memory function. Second, the same spectral
density fixes both the branch cut of Section~\ref{sec:propagator} and the
long-distance $1/r^3$ residue of Section~\ref{sec:longdistance} --- they
are one prediction, not two. This structure agrees with the explicit
integro-differential brane perturbation equation of Mukohyama
\cite{Mukohyama2}.

\subsection{Nonlinear completion}

At the nonlinear level there is generally no vector-space kernel and
hence no quotient by a subspace. Instead:

\begin{definition}[Operational equivalence of bulk data]
\label{def:nonlinearEquiv}
Two admissible gauge-reduced bulk data sets are equivalent when, for the
same retained brane initial data, they induce the same complete brane
record for every allowed brane protocol. The minimal predictive
completion is the set of these equivalence classes.
\end{definition}

The tensor $E_{\mu\nu}$ is the \emph{local manifestation} of the missing
bulk response, not a substitute for it. Its divergence equation
\eqref{eq:bianchiE} constrains it but does not evolve it; supplying
$E_{\mu\nu}$ instantaneously does not close the system, because its
subsequent evolution requires the bulk problem. The nonlinear completion
is therefore neither the linear quotient \eqref{eq:Edyn} nor, in general,
a finite set of four-dimensional functions.

\section{What the source model fixes}
\label{sec:fixed}

Conditional on RS2, the source action fixes the following:
\begin{enumerate}[leftmargin=*]
\item the tensor form and numerical coefficients of $\Pi_{\mu\nu}$,
\eqref{eq:Pi};
\item the sign of the high-density correction for a positive-tension
brane;
\item the $\rho^2$ and $a^{-4}$ scalings in the homogeneous sector;
\item the massless graviton pole plus a positive KK continuum in linear
RS2;
\item the leading $1/r^3$ potential correction and the ratio of the two
metric potentials;
\item the cross-regime relation \eqref{eq:crossregime};
\item covariant conservation of the total projected response.
\end{enumerate}

Everything \emph{not} fixed --- the choice of RS2 itself, the value of
$\ell$, the bulk state and $\mu$, the boundary condition, possible
non-minimal operators, and the ultraviolet completion --- is collected
with the corresponding failure modes in Section~\ref{sec:failure}. The
cross-regime relation removes $\ell$ from one comparison; it does not
remove the physical assumptions defining the source model.

\section{The strongest conventional four-dimensional description}
\label{sec:fourdim}

\subsection{Homogeneous sector}

Define
\begin{equation}
\label{eq:Teff}
T^{\mathrm{eff}}_{\mu\nu}
=T_{\mu\nu}
+\frac{\kfive^4}{\kfour^2}\Pi_{\mu\nu}
-\frac{1}{\kfour^2}E_{\mu\nu}.
\end{equation}
The projected equation \eqref{eq:SMS} becomes ordinary four-dimensional
Einstein gravity sourced by $T^{\mathrm{eff}}_{\mu\nu}$. In the
homogeneous sector,
\begin{equation}
\label{eq:rhoeff}
\rho_{\mathrm{eff}}
=\rho+\frac{\rho^2}{2\lambda}
+\frac{3\mu}{\kfour^2a^4},
\qquad
p_{\mathrm{eff}}
=p+\frac{\rho(\rho+2p)}{2\lambda}
+\frac{\mu}{\kfour^2a^4}.
\end{equation}
This four-dimensional fluid description reproduces the background
expansion \emph{exactly}, and may perfectly well be regarded as
fundamental rather than induced. Background cosmological data alone
cannot decide between the two readings.

\subsection{\texorpdfstring{Linear sector, and the CFT realization of the $1/r^3$ residue}{Linear sector, and the CFT realization of the 1/r\textthreesuperior{} residue}}

The quadratic action of Section~\ref{sec:linear} is already a
four-dimensional \emph{nonlocal} action once $\mathcal G_R$ is taken as
its primitive kernel. That such nonlocal braneworld actions with explicit
KK form factors can be written down constructively is shown, in the
two-brane setting, by Barvinsky \emph{et al.} \cite{Barvinsky}; we invoke
that work only as a constructive control on the \emph{form} of the
eliminated action, not as an identity for RS2, whose single-brane
spectrum is a continuum rather than a discrete tower. It contains the same massless pole, the same
continuum branch cut, the same causal response, and the same weak-field
coefficients as the five-dimensional description.

The four-dimensional alternative is not hypothetical. The stress-tensor
vacuum-polarization loop of an appropriate four-dimensional conformal
sector corrects the graviton propagator with the \emph{same} $1/r^3$
coefficient as the classical RS2 bulk calculation \cite{DuffLiu,APR,PerezVictoria}.
Thus even the most characteristic weak-field residue of the model has a
conventional four-dimensional realization. This is the concrete form
taken, inside RS2, by the general obstruction proved in
Section~\ref{sec:opeq}.

The discrimination supplied by the model is therefore
\begin{equation}
\label{eq:discrim}
\text{minimal RS2}
\quad\text{versus}\quad
\text{4D GR or another \emph{independently restricted} model class},
\end{equation}
and \emph{not}
\begin{equation}
\label{eq:notdiscrim}
\text{source projection}
\quad\text{versus}\quad
\text{every possible brane-only theory}.
\end{equation}

\section{The operational-equivalence obstruction}
\label{sec:opeq}

Sections~\ref{sec:witness}--\ref{sec:completion} establish nonclosure of
the \emph{instantaneous} brane state. That statement must be sharply
distinguished from the far stronger claim that no autonomous description
using only brane observables can exist. The stronger claim is
\emph{false} when arbitrary history dependence is allowed.

\subsection{Classical pushforward lemma}

Let $\mathfrak E_T$ be the class of all experiments that prepare,
intervene on, and read only Tier-1 (brane) variables. For each protocol
$e\in\mathfrak E_T$, a fully specified source theory defines a
probability measure $\mu^S_e$ on its admissible source histories
$\Gamma^e_S$. Let
\begin{equation}
\label{eq:Re}
R_e:\Gamma^e_S\longrightarrow\Gamma^e_T
\end{equation}
be the measurable map returning the complete Tier-1 record.

\begin{lemma}[Pushforward representation of Tier-1 operational laws]
\label{lem:pushforward}
Define
\begin{equation}
\label{eq:nu}
\boxed{\;
\nu_e=(R_e)_*\mu^S_e,
\qquad e\in\mathfrak E_T .
\;}
\end{equation}
Then $M^\star_T=\{\nu_e:e\in\mathfrak E_T\}$ is a theory whose sample
space contains only Tier-1 histories and whose prediction for every
Tier-1 experiment coincides exactly with the source prediction:
\begin{equation}
\label{eq:TVzero}
d_{\mathrm{TV}}\bigl((R_e)_*\mu^S_e,\ \nu^\star_e\bigr)=0
\qquad\text{for every }e\in\mathfrak E_T ,
\end{equation}
where $\nu^\star_e:=\nu_e$ regarded as the surrogate theory's law.
\end{lemma}

\begin{proof}
By the definition of pushforward, for every bounded measurable statistic
$F:\Gamma^e_T\to\mathbb R$,
\begin{equation}
\label{eq:changeofvars}
\int_{\Gamma^e_T}F(\gamma)\,\nu_e(d\gamma)
=\int_{\Gamma^e_S}F(R_e\omega)\,\mu^S_e(d\omega).
\end{equation}
Taking $F$ to be the indicator of an arbitrary measurable event proves
equality of the complete record distributions, hence \eqref{eq:TVzero}.
Taking $F$ to encode an estimator, a randomized decision rule, a stopping
rule, a correlation function, or a loss function proves equality of its
distribution as well. Adaptive interventions are already covered, because
each complete adaptive protocol \emph{is} a member of $\mathfrak E_T$ and
its record map is measurable. Marginalization, random mixtures of
protocols, and causal consistency are inherited because pushforward
commutes with measurable composition and with marginalization.
\end{proof}

\begin{remark}[Status of Lemma~\ref{lem:pushforward}]
\label{rem:pushstatus}
The equality \eqref{eq:TVzero} is mathematically immediate from the
definition of pushforward; nothing in the proof is deep, and no claim of
mathematical originality is made for it. Its significance here is
\emph{interpretive}. It establishes the exact boundary on what brane-only
observations can discriminate, and hence fixes the strongest conclusion
that the positive results of
Sections~\ref{sec:witness}--\ref{sec:crossregime} can support. The
substantive work lies in identifying which restrictions on the alternative
class \emph{do} permit discrimination
(Section~\ref{sec:boundary}), and in exhibiting --- through the
four-dimensional CFT description of Section~\ref{sec:fourdim} --- that the
obstruction is realized concretely in RS2 and not merely as a formal
possibility.
\end{remark}

For standard Borel state spaces the finite-time laws disintegrate into
causal history kernels,
\begin{equation}
\label{eq:disintegration}
\nu(dx_{0:n}\mid u_{0:n-1})
=K_0(dx_0)\prod_{j=0}^{n-1}
K_{j+1}\bigl(dx_{j+1}\mid x_{0:j},u_{0:j}\bigr),
\end{equation}
an autonomous stochastic law on Tier-1 histories and controls. It need
not be Markovian, finite-dimensional, local, simple, or state-independent.
Those are \emph{explanatory} restrictions, not consequences of autonomy
in the unrestricted operational sense.

\subsection{Quantum and field-theoretic forms}

In quantum theory, equality of one-time reduced states, or of an ordinary
in--out generating functional, is insufficient: interventions can expose
temporal correlations. The exact Tier-1 object is the \emph{process
tensor} \cite{Pollock}. If $\mathcal A^{(j)}_{x_j}$ are Tier-1
instruments and $\mathcal U_{j+1:j}$ the joint source propagator, then
\begin{equation}
\label{eq:processprob}
\Pr(\mathbf x\mid\mathbf A)
=\Tr_{TE}\Bigl[
\mathcal A^{(n)}_{x_n}\,\mathcal U_{n:n-1}\cdots
\mathcal A^{(0)}_{x_0}(\rho_{TE})\Bigr].
\end{equation}
Contracting every eliminated-sector index defines a positive, causally
normalized multilinear functional $\Upsilon^T_{n:0}$ with Tier-1 legs
only, such that
\begin{equation}
\label{eq:processtensor}
\Pr(\mathbf x\mid\mathbf A)
=\Upsilon^T_{n:0}\bigl[
\mathcal A^{(n)}_{x_n},\mathcal A^{(n-1)}_{x_{n-1}},\ldots,
\mathcal A^{(0)}_{x_0}\bigr].
\end{equation}

The field-theoretic counterpart is the Feynman--Vernon closed-time-path
influence functional \cite{FeynmanVernon}. With
$S[\phi,\chi]=S_T[\phi]+S_E[\chi]+S_{\mathrm{int}}[\phi,\chi]$ and a
factorized initial environment state,
\begin{equation}
\label{eq:IF}
e^{iS_{\mathrm{IF}}[\phi_+,\phi_-]}
=\int_{\chi_+(t_f)=\chi_-(t_f)}
\mathcal D\chi_+\,\mathcal D\chi_-\;
\rho_E(\chi^+_i,\chi^-_i)\,
e^{iS_E[\chi_+]+iS_{\mathrm{int}}[\phi_+,\chi_+]
-iS_E[\chi_-]-iS_{\mathrm{int}}[\phi_-,\chi_-]},
\end{equation}
\begin{equation}
\label{eq:ZT}
Z_T[J_+,J_-]
=\int\mathcal D\phi_+\,\mathcal D\phi_-\;
\rho_T(\phi^+_i,\phi^-_i)\,
e^{iS_T[\phi_+]-iS_T[\phi_-]
+iS_{\mathrm{IF}}[\phi_+,\phi_-]
+iJ_+\phi_+-iJ_-\phi_-} .
\end{equation}
The retained action is counted exactly once. A \emph{correlated} initial
state replaces $\rho_E$ by the corresponding joint boundary functional;
the process-tensor construction \eqref{eq:processtensor} covers that case
with no factorization assumption. The resulting closed-time-path
functional contains retarded response, fluctuations, and multi-time
Tier-1 probabilities, and may be complex, nonlocal, and state-dependent
while remaining a complete Tier-1 description.

\subsection{The exact boundary of the conclusion}
\label{sec:boundary}

The witness of Section~\ref{sec:witness} proves
\begin{equation}
\label{eq:noflow}
p\circ\Phi_t\ne U_t\circ p
\end{equation}
for every deterministic flow $U_t$ on the specified instantaneous brane
state, and --- for any preparation class that resolves the two witness
states, so that the conditional future law differs within a readout fiber
--- it also defeats universal Markov lumpability. It does
\textbf{not} defeat the pushed-forward law \eqref{eq:nu} on complete
brane histories.

Accordingly, source projection can be empirically distinguished only
relative to a \emph{restricted} alternative class $\mathfrak M^R_T$. If
$\Op(M)$ denotes the complete family of protocol-conditioned Tier-1 laws
of a model, the necessary and sufficient condition that the restricted
class contain no exactly operationally equivalent model is
\begin{equation}
\label{eq:identifiability}
\boxed{\;
\Op(M_S,p)\notin\bigl\{\Op(M):M\in\mathfrak M^R_T\bigr\}.
\;}
\end{equation}

\begin{remark}[Identifiability is not a finite-data testing theorem]
\label{rem:finitedata}
Condition \eqref{eq:identifiability} is exact identifiability. Uniform
statistical discrimination from a \emph{composite} alternative class
additionally requires separation of the source law from the
\emph{closure} of that class in an operational metric appropriate to the
allowed protocols --- total variation, or a relative-entropy rate. If the
source law lies outside the class but can be approximated arbitrarily
closely by its members, exact equivalence is absent while finite-data
discrimination may nonetheless be arbitrarily difficult.
\end{remark}

Useful independent restrictions include: a fixed finite set of free or
weakly coupled instantaneous fields at a stated perturbative order; a
Markov or finite-memory condition; fixed finite-pole local quadratic
dynamics; a finite EFT basis with coefficients fixed before the test; a
prescribed preparation/boundary-state class. \emph{Locality alone, or
finitely many interacting field species alone, is insufficient}, because
multiparticle states generate branch cuts. If arbitrary kernels,
continua, initial correlations, and complete-history states are admitted,
the pushforward theory \emph{belongs} to the alternative class and
discrimination is impossible.

\section{Main theorem}
\label{sec:main}

\begin{theorem}[RS2 source projection and its exact discrimination boundary]
\label{thm:main}
Consider the two-sided, $Z_2$-symmetric, positive-tension RS2 model
defined by \eqref{eq:action}, with pure five-dimensional Einstein gravity
and a negative cosmological constant in the bulk. Tune
$\Lfive=-6/\ell^2$ and $\lambda=6/(\kfive^2\ell)$, so that $\Lfour=0$,
and let the instantaneous readout be Definition~\ref{def:readout}.
Assume the full source domain contains smooth finite-energy TT bulk
initial data. For the stationary weak-field coefficients, impose the
regular retarded/no-incoming bulk condition. Then:
\begin{enumerate}[leftmargin=*]
\item the readout has physically distinct fibers on which future induced
brane observables differ, so no deterministic evolution on that
instantaneous readout is valid for all admissible bulk initial data;
\item the exact projected tensor equation contains the model-specific
correction
$\mathcal C^{\RS}_{\mu\nu}=\kfive^4\Pi_{\mu\nu}-E_{\mu\nu}$
with $\Pi_{\mu\nu}$ as in \eqref{eq:Pi}, and predictive closure requires
either bulk response data or an equivalent causal, generally nonlocal
brane history law;
\item in the homogeneous and long-distance weak-field regimes the
observable coefficients obey
$A_\Psi=\tfrac12A_\Phi$ and $B=\tfrac{\kfour^4}{24}A_\Phi$;
\item these statements distinguish minimal RS2 from any independently
specified Tier-1 model class that cannot realize the correlated quadratic
term, slip ratio, and continuum response --- but no experiment restricted
to brane observables distinguishes the five-dimensional source model from
its \emph{unrestricted} exact four-dimensional pushforward theory.
\end{enumerate}
\end{theorem}

\begin{proof}
\textbf{Part 1.} Take zero bulk perturbation and a nonzero, normalizable
TT wave packet initially supported away from the brane, with the
perturbation and its first time derivative vanishing in a brane
neighborhood. The two induced Cauchy data sets coincide
(Proposition~\ref{prop:witness}). Their five-dimensional Weyl curvatures
do not, so they are not related by a diffeomorphism. Hyperbolic finite
propagation permits the packet to reach the brane and change the induced
perturbation. Hence one readout fiber is not preserved by source
evolution, contradicting the defining condition for a deterministic
quotient flow (Corollary~\ref{cor:nonproj}).

\textbf{Part 2.} The Israel equation \eqref{eq:israel} with $Z_2$
symmetry gives \eqref{eq:Kfixed}. The contracted Gauss identity separates
the bulk Riemann tensor into its Ricci part, fixed by $\Lfive$ via
\eqref{eq:bulkEE}, and its Weyl part $E_{\mu\nu}$. The terms quadratic in
$K_{\mu\nu}$ expand into a constant, a term linear in $T_{\mu\nu}$, and
$\kfive^4\Pi_{\mu\nu}$, as computed line by line in
\eqref{eq:Qexpanded}--\eqref{eq:Qsorted} with $\Pi_{\mu\nu}$ given by
\eqref{eq:Pi},
giving
\begin{equation}
G_{\mu\nu}
=-\tfrac12\Bigl(\Lfive+\tfrac{\kfive^4\lambda^2}{6}\Bigr)q_{\mu\nu}
+\frac{\kfive^4\lambda}{6}T_{\mu\nu}
+\kfive^4\Pi_{\mu\nu}-E_{\mu\nu},
\end{equation}
which proves \eqref{eq:SMS} and \eqref{eq:constants}. The Bianchi
identity fixes only $\nabla^\mu E_{\mu\nu}$, \eqref{eq:bianchiE}. Part 1
supplies two solutions with the same instantaneous brane data and
different later $E_{\mu\nu}$; the remaining information must therefore be
supplied as bulk initial/boundary response data or, after exact
elimination, as a causal functional of the brane history and the chosen
bulk state (Section~\ref{sec:completion}).

\textbf{Part 3.} Perfect-fluid substitution gives $\Pi_{00}=\rho^2/12$
(Lemma~\ref{lem:PiFRW}), while FRW symmetry gives $-E_{00}=3\mu/a^4$
\eqref{eq:EFRW}. The $00$ equation therefore gives $B=\kfive^4/36$. RS
tuning with $\kfour^2=\kfive^2/\ell$ gives $B=\kfour^4\ell^2/36$.
Independently, the regular brane-to-brane propagator gives
$A_\Phi=2\ell^2/3$ and $A_\Psi=\ell^2/3$ at $r\gg\ell$
\eqref{eq:potentials}. Eliminating $\ell^2$ proves
$A_\Psi=A_\Phi/2$ and $B=\kfour^4A_\Phi/24$
(Theorem~\ref{thm:crossregime}).

\textbf{Part 4.} A restricted theory that cannot realize one of these
fixed relations is falsified when all assumptions and coefficients are
fixed before observing the test data. Conversely, for every brane
protocol, push the RS2 source-history measure forward through the
complete brane readout. Lemma~\ref{lem:pushforward} shows the resulting
four-dimensional history theory has exactly the same record distribution
for that protocol, including adaptive interventions; their total-variation
distance vanishes for every brane experiment. This proves both the
positive antecedent and the negative ontological-discrimination
statement.
\end{proof}

Theorem~\ref{thm:main} establishes a concrete, quantitative
source-projection realization together with its exact discrimination
boundary. A stronger conclusion requires an independently justified
restriction on admissible Tier-1 dynamics --- fixed finite-pole local
quadratic dynamics, prescribed finite perturbative content, or finite
memory.

\section{Comparative examples and limits of interpretation}
\label{sec:controls}

\subsection{An exactly soluble projection: two coupled oscillators}
\label{sec:oscillator}

The following systems are scientific comparisons that clarify the meaning
and limits of the RS2 result. An ordinary Hamiltonian model verifies, without gravitational
complications, every step from fiber nonprojectability through memory to
an exact Tier-1 surrogate, and simultaneously exhibits the logical
boundary that Section~\ref{sec:opeq} proves in general. The projected
equation obtained below is the elementary exactly soluble instance of the
Mori--Zwanzig projection-operator structure \cite{MoriZwanzig}: a memory
kernel, an initial-slip term, and a retained susceptibility. Let
\begin{equation}
\label{eq:Hosc}
H_f=\frac{p_x^2}{2m}+\tfrac12 m\omega_0^2x^2
+\frac{p_y^2}{2M}
+\tfrac12 M\Omega^2\Bigl(y-\frac{cx}{M\Omega^2}\Bigr)^2
-f(t)\,x ,
\end{equation}
and retain $X=(x,v)^T$ while eliminating $Z=(y,w)^T$. With
$\delta\equiv c^2/(mM\Omega^2)$, the first-order generator has blocks
\begin{equation}
\label{eq:blocks}
A=\begin{pmatrix}0&1\\-(\omega_0^2+\delta)&0\end{pmatrix},
\quad
B=\begin{pmatrix}0&0\\c/m&0\end{pmatrix},
\quad
C=\begin{pmatrix}0&0\\c/M&0\end{pmatrix},
\quad
D=\begin{pmatrix}0&1\\-\Omega^2&0\end{pmatrix},
\end{equation}
with
$e^{tD}=\begin{pmatrix}\cos\Omega t&\sin\Omega t/\Omega\\
-\Omega\sin\Omega t&\cos\Omega t\end{pmatrix}$.
Variation of constants evaluates \emph{every} projected term:
\begin{equation}
\label{eq:oscproj}
\dot X(t)=AX(t)+Be^{tD}Z_0
+\int_0^t\mathcal K(t-s)X(s)\,ds
+\binom{0}{f(t)/m},
\end{equation}
\begin{equation}
\label{eq:oscK}
\boxed{\;
\mathcal K(t)=Be^{tD}C
=\begin{pmatrix}
0&0\\[2pt]
\dfrac{c^2}{mM\Omega}\sin\Omega t&0
\end{pmatrix}.
\;}
\end{equation}
Equivalently, after one integration by parts,
\begin{equation}
\label{eq:oscGLE}
\boxed{\;
\ddot x+\omega_0^2x
+\int_0^t\delta\cos\bigl[\Omega(t-s)\bigr]\dot x(s)\,ds
=\frac{f(t)}{m}+\zeta(t),
\;}
\end{equation}
\begin{equation}
\label{eq:oscnoise}
\zeta(t)=\frac{c}{m}
\Bigl[\Bigl(y_0-\frac{cx_0}{M\Omega^2}\Bigr)\cos\Omega t
+\frac{w_0}{\Omega}\sin\Omega t\Bigr].
\end{equation}
Two states with equal $(x_0,v_0)$ and different $(y_0,w_0)$ have
different later trajectories --- the instantaneous readout is
nonprojectable, exactly as in Proposition~\ref{prop:witness}, with
$\zeta$ playing the role of $\eta_{z_0}$ in \eqref{eq:memoryeq}. The
smallest generic Markov completion carries the two hidden coordinates,
exactly as $E^{\RS}_{\mathrm{dyn}}$ does. The exact retarded
susceptibility is
\begin{equation}
\label{eq:chi}
\boxed{\;
\chi^R_{xx}(\omega)
=\frac{\Omega^2-z^2}
{m\bigl[(\omega_0^2-z^2)(\Omega^2-z^2)-\delta z^2\bigr]},
\qquad z=\omega+i0^+ .
\;}
\end{equation}

And yet eliminating $y$ also gives a \emph{local, Tier-1-only},
higher-derivative equation,
\begin{equation}
\label{eq:oscfourth}
\boxed{\;
\bigl[(\partial_t^2+\omega_0^2)(\partial_t^2+\Omega^2)
+\delta\,\partial_t^2\bigr]x
=\frac1m(\partial_t^2+\Omega^2)f .
\;}
\end{equation}
Its extra initial derivatives encode precisely the two missing constants.
This exactly soluble example exposes the logical boundary: nonprojectability
of $(x,\dot x)$ proves the need for completion \emph{relative to that
state choice}, and proves nothing about the nonexistence of a
higher-derivative or history-level theory written in $x$ alone. This is
the finite-dimensional shadow of Lemma~\ref{lem:pushforward}.

\subsection{Einstein--Cartan and Aharonov--Bohm comparisons}
\label{sec:countercontrols}

These two examples distinguish a genuine propagating source degree of
freedom from an auxiliary variable and from an incomplete choice of
observable state: Einstein--Cartan--Dirac torsion is auxiliary and
introduces no same-readout fiber, whereas the Aharonov--Bohm phase is
already determined once the observable state is chosen completely.

\paragraph{Einstein--Cartan--Dirac.}
Minimal Einstein--Cartan--Dirac theory fails as the required
\emph{propagating source-projection antecedent} --- not as a theory ---
because its torsion is \emph{auxiliary}: it introduces no independent
propagating degree of freedom, and hence no same-readout fiber of the
kind Proposition~\ref{prop:witness} exhibits. The connection equation
determines the contorsion algebraically from the axial current, and
substitution gives the four-fermion term \cite{Diakonov}
\begin{equation}
\label{eq:ECD}
S_{\mathrm{eff}}
=S_{\mathrm{EH}}+S_{\mathrm{Dirac}}[\mathring\omega]
+\frac{3\kappa}{16}\int d^4x\;e\,J_{5a}J_5^a,
\end{equation}
up to signature and current conventions. A torsion-free theory containing
this operator has \emph{identical} retained observables. There is no
independent propagating fiber analogous to a remote RS2 wave packet ---
which is exactly what RS2 supplies and Einstein--Cartan does not.

\paragraph{Aharonov--Bohm.}
The Aharonov--Bohm effect \cite{AharonovBohm} remains a clean control for
readout incompleteness: equal local curvature can coexist with different
gauge-invariant holonomy and different interference phases. But holonomy
is already a standard observable of four-dimensional gauge theory. The
effect shows that a \emph{curvature-only} state is incomplete; it does
not show that a deeper source ontology has been observed. It is the
correct model for ``the readout you chose was too small,'' not for
``there is a hidden dimension.''

\section{Empirical discrimination}
\label{sec:empirical}

\subsection{A pre-registered model test}

Fit the three phenomenological coefficients of
\eqref{eq:phenomB}--\eqref{eq:phenomA} \emph{without} using RS2 in the
fits, and form the null combinations
\begin{equation}
\label{eq:nulls}
\boxed{\;
D_1=A_\Phi-2A_\Psi,
\qquad
D_2=24B-\kfour^4A_\Phi .
\;}
\end{equation}
Minimal RS2 predicts $D_1=D_2=0$ with $A_\Phi,A_\Psi,B\ge0$.

\begin{remark}[These are consistency tests, not a detection criterion]
Pure four-dimensional GR has $A_\Phi=A_\Psi=B=0$ and passes both nulls
trivially. Discrimination from GR requires at least one \emph{statistically
significant positive} coefficient together with consistency of the
others. The coefficients must be fit from independent observables, with
the functional forms fixed before examining residuals.
\end{remark}

In the nonrelativistic sector,
\begin{equation}
\label{eq:accel_frac}
\frac{\Delta|\mathbf a|}{|\mathbf a|}
=\frac{3A_\Phi}{r^2}+O(r^{-4}),
\end{equation}
whereas relativistic propagation probes $\Phi+\Psi$ and therefore the
different combination $(A_\Phi+A_\Psi)/2=3A_\Phi/4$. The two probes are
not degenerate, which is what prevents a single rescaling of $G_4M$ from
absorbing the signal.

Short-range torsion-balance experiments already test the inverse-square
law at submillimetre separations \cite{Lee}. Such data should be fitted
to the \emph{full} RS2 KK spectral potential built from \eqref{eq:um},
convolved with the extended attractor and detector geometry, rather than
by translating a Yukawa exclusion curve. The truncated $1/r^3$ potential
may be used only after the fitted parameter range verifies $r\gg\ell$;
near $r\sim\ell$ it is uncontrolled. Primordial-abundance and expansion
data constrain the nonstandard high-density expansion and the separate
radiation-like term \cite{Ichiki}. A combined analysis must fit $B$ and
$\mu$ \emph{separately}: during radiation domination $\rho^2\propto
a^{-8}$ whereas dark radiation scales as $a^{-4}$, so a single-epoch
expansion residual can mix their amplitudes, but multi-epoch data and
equation-of-state transitions separate them.

This is a falsification test, not a discovery forecast. Existing
short-distance bounds already make the corresponding high-energy
cosmological term small at late epochs, and a simultaneous precision
measurement of both static potentials at microscopic distances is
technically difficult. At present the cross-regime relation is more
useful as a consistency condition on any \emph{claimed} anomaly than as a
near-term detection channel.

\subsection{Spectral test}

A time-dependent brane source couples to a massless pole plus a
continuum. A response measurement with sufficient bandwidth could compare
a finite-pole local model against the positive KK spectral measure
\begin{equation}
\label{eq:rhoKK}
\rho_{\KK}(m)=u_m(0)^2\ge0,
\qquad m\in[0,\infty).
\end{equation}
A resolved threshold continuum, or the corresponding power-law late-time
tail, would exclude a \emph{pre-specified finite-pole free or tree-level
linear} alternative. It would \emph{not} exclude an interacting
four-dimensional conformal sector or continuum model, which has the same
analytic structure (Remark~\ref{rem:cut}).

\subsection{What no experiment of this kind can decide}

No statistic constructed solely from induced brane fields can decide
between RS2 and the exact Tier-1 process functional built from RS2. This
is not a limitation of experimental sensitivity; it follows from equality
of probability measures, \eqref{eq:TVzero}. An experiment that couples
directly to a bulk coordinate could distinguish the source fiber, but it
would thereby enlarge the observable algebra and would no longer be
confined to the readout used in the theorem.

\section{Scope and regime of validity}
\label{sec:failure}

The results derived above apply to the following physical domain and
assumptions.

\begin{enumerate}[leftmargin=*]
\item \textbf{Bulk-state domain.} The same-readout witness requires that
admissible initial data include normalizable bulk radiation. If the model
definition fixes a unique stationary vacuum and excludes all homogeneous
incoming data, the witness lies outside the restricted domain. The forced
KK response and the effective-action equivalence remain, but the claimed
instantaneous fiber must be reassessed for that domain.

\item \textbf{Minimal action.} A brane Einstein term, a Gauss--Bonnet
term, additional bulk fields, finite brane thickness, a second brane, or
broken $Z_2$ symmetry changes the junction condition, the propagator, and
the coefficient relation. In such theories $D_1=D_2=0$ need not hold.

\item \textbf{Positive-tension branch.} The sign and stability statements
use $\lambda>0$. The negative-tension branch can reverse kinetic signs or
the effective Newton coupling; it is not a harmless sign convention.

\item \textbf{Homogeneous sector.} Equation \eqref{eq:friedmann} assumes
an FRW brane, conserved brane matter, and no bulk matter flux. With bulk
exchange, both $\dot\rho+3H(\rho+p)=0$ and $\mathcal U\propto a^{-4}$ are
modified.

\item \textbf{Weak-field sector.} The coefficients $A_\Phi,A_\Psi$ assume
a static weak source, $r\gg\ell$, and the regular retarded/no-incoming
bulk condition. They are unreliable in the strong-field region and at
$r\lesssim\ell$.

\item \textbf{The state parameter $\mu$.} Neither the sign nor the
magnitude of $\mu$ is fixed by the projected brane equation. Positivity
follows only after imposing a positive-mass regular AdS--Schwarzschild
bulk. Treating dark radiation as a universal fixed prediction would
overstate the model.

\item \textbf{Observable versus decomposition.} The induced metric and
matter records are observable; the \emph{split}
$\mathcal C^{\RS}=\kfive^4\Pi-E$ is source-model dependent. A
four-dimensional theory may reorganize the same total response into
different fields or a nonlocal kernel.

\item \textbf{Unrestricted alternatives.} Nonlocality, colored noise, a
branch cut, or preparation dependence does not by itself refute a Tier-1
theory. Each is permitted in the exact pushed-forward history law
\eqref{eq:nu}. A contrary conclusion requires an independently imposed
restriction --- finite-pole local quadratic dynamics, fixed perturbative
content, or Markovianity. Locality, or finitely many interacting fields,
is \emph{not} enough.

\item \textbf{Flat-brane tuning.} The cross-regime relation was derived
at $\Lfour=0$. A detuned or de Sitter brane changes the background
junction relation and can modify the weak-field coefficients. Applying
the flat-brane relation to a weakly curved universe is an approximation
requiring $H\ell\ll1$ and negligible curvature/detuning effects; it is not
a complete late-time cosmological fit.
\end{enumerate}

Two further parameters are simply undetermined rather than failure modes.
The numerical value of $\ell$ --- equivalently $\lambda$ or $\kfive$ once
$G_4$ is fixed --- is not predicted; only bounds on it follow from data.
And the theory is a five-dimensional effective field theory, so nothing
here constrains its ultraviolet completion above the corresponding cutoff.

\paragraph{Decoupling limits.}
These are consistent. As $\ell\to0$ at fixed $G_4$: $A_\Phi,A_\Psi\to0$,
$\lambda\to\infty$, and $B\to0$. As $\rho/\lambda\to0$,
$\mu/(a^4H^2)\to0$, and $r/\ell\to\infty$, ordinary four-dimensional GR
is recovered in every sector.

\section{Conclusion}
\label{sec:conclusion}

Minimal positive-tension RS2 contains a complete source-projection
antecedent. In that model, the projection from gauge-reduced bulk gravity
to the instantaneous brane state removes genuine gravitational degrees of
freedom: a remote transverse-traceless bulk wave, whose existence is
established by an explicit KK spectral construction, supplies a
same-readout, different-future pair, so no deterministic evolution law
exists on that instantaneous state over the full source domain.
Gauss--Codazzi projection and Israel matching \emph{derive}, rather than
posit,
\begin{equation}
\mathcal C^{\RS}_{\mu\nu}=\kfive^4\Pi_{\mu\nu}-E_{\mu\nu}.
\end{equation}
Its local nonlinear part gives the exact positive $\rho^2$ correction;
its bulk Weyl part carries the completion data and the retarded KK memory.
The same source scale $\ell$ fixes the leading weak-field potentials,
producing the parameter-eliminated relation
\begin{equation}
A_\Psi=A_\Phi/2,
\qquad
B=\kfour^4A_\Phi/24 ,
\end{equation}
a genuine, falsifiable residue linking two otherwise independent
observational regimes. This is the strongest quantitative result obtained
here.

The strongest possible \emph{interpretation} is nevertheless bounded by
Lemma~\ref{lem:pushforward}: every source model induces, by pushforward, a
brane-only path law with identical brane statistics, and RS2 makes that
obstruction concrete through its four-dimensional CFT description, which
reproduces the same $1/r^3$ residue. The result is therefore positive at
the level of model-relative source projection and nonclosure, and negative
at the level of unrestricted ontological discrimination. Absent an
independently motivated restriction on the admissible brane-only theory
class, no brane observation can establish that the source representation
is ontologically fundamental; and establishing that nature realizes this
antecedent at all requires empirical selection of RS2, or of another source
model with comparably specific residues. The conditions under which the
positive results hold are stated once, in Section~\ref{sec:failure}.

\begin{thebibliography}{99}

\bibitem{RS2}
L.~Randall and R.~Sundrum,
``An Alternative to Compactification,''
\emph{Phys.\ Rev.\ Lett.} \textbf{83}, 4690 (1999),
arXiv:hep-th/9906064.

\bibitem{SMS}
T.~Shiromizu, K.~Maeda, and M.~Sasaki,
``The Einstein Equations on the 3-Brane World,''
\emph{Phys.\ Rev.\ D} \textbf{62}, 024012 (2000),
arXiv:gr-qc/9910076.

\bibitem{Israel}
W.~Israel,
``Singular Hypersurfaces and Thin Shells in General Relativity,''
\emph{Nuovo Cim.\ B} \textbf{44}, 1 (1966),
doi:10.1007/BF02710419.

\bibitem{BDEL}
P.~Binetruy, C.~Deffayet, U.~Ellwanger, and D.~Langlois,
``Brane Cosmological Evolution in a Bulk with Cosmological Constant,''
\emph{Phys.\ Lett.\ B} \textbf{477}, 285 (2000),
arXiv:hep-th/9910219.

\bibitem{Mukohyama1}
S.~Mukohyama,
``Brane-World Solutions, Standard Cosmology, and Dark Radiation,''
\emph{Phys.\ Lett.\ B} \textbf{473}, 241 (2000),
arXiv:hep-th/9911165.

\bibitem{GarrigaTanaka}
J.~Garriga and T.~Tanaka,
``Gravity in the Randall--Sundrum Brane World,''
\emph{Phys.\ Rev.\ Lett.} \textbf{84}, 2778 (2000),
arXiv:hep-th/9911055.

\bibitem{GKR}
S.~B.~Giddings, E.~Katz, and L.~Randall,
``Linearized Gravity in Brane Backgrounds,''
\emph{JHEP} \textbf{03}, 023 (2000),
arXiv:hep-th/0002091.

\bibitem{Mukohyama2}
S.~Mukohyama,
``Integro-Differential Equation for Brane-World Cosmological
Perturbations,''
\emph{Phys.\ Rev.\ D} \textbf{64}, 064006 (2001),
arXiv:hep-th/0104185.

\bibitem{DuffLiu}
M.~J.~Duff and J.~T.~Liu,
``Complementarity of the Maldacena and Randall--Sundrum Pictures,''
\emph{Phys.\ Rev.\ Lett.} \textbf{85}, 2052 (2000),
arXiv:hep-th/0003237.

\bibitem{APR}
N.~Arkani-Hamed, M.~Porrati, and L.~Randall,
``Holography and Phenomenology,''
\emph{JHEP} \textbf{08}, 017 (2001),
arXiv:hep-th/0012148.

\bibitem{Barvinsky}
A.~O.~Barvinsky, A.~Yu.~Kamenshchik, A.~Rathke, and C.~Kiefer,
``Nonlocal Braneworld Action: An Alternative to Kaluza--Klein
Description,''
\emph{Phys.\ Rev.\ D} \textbf{67}, 023513 (2003),
arXiv:hep-th/0206188.

\bibitem{FeynmanVernon}
R.~P.~Feynman and F.~L.~Vernon, Jr.,
``The Theory of a General Quantum System Interacting with a Linear
Dissipative System,''
\emph{Ann.\ Phys.} \textbf{24}, 118 (1963),
doi:10.1016/0003-4916(63)90068-X.

\bibitem{MoriZwanzig}
H.~Mori,
``Transport, Collective Motion, and Brownian Motion,''
\emph{Prog.\ Theor.\ Phys.} \textbf{33}, 423 (1965),
doi:10.1143/PTP.33.423;
R.~Zwanzig,
``Nonlinear Generalized Langevin Equations,''
\emph{J.\ Stat.\ Phys.} \textbf{9}, 215 (1973),
doi:10.1007/BF01008729.

\bibitem{Pollock}
F.~A.~Pollock, C.~Rodr\'iguez-Rosario, T.~Frauenheim, M.~Paternostro,
and K.~Modi,
``Non-Markovian Quantum Processes: Complete Framework and Efficient
Characterisation,''
\emph{Phys.\ Rev.\ A} \textbf{97}, 012127 (2018),
arXiv:1512.00589.

\bibitem{Diakonov}
D.~Diakonov, A.~G.~Tumanov, and A.~A.~Vladimirov,
``Low-Energy General Relativity with Torsion: A Systematic Derivative
Expansion,''
\emph{Phys.\ Rev.\ D} \textbf{84}, 124042 (2011),
arXiv:1104.2432.

\bibitem{Lee}
J.~G.~Lee, E.~G.~Adelberger, T.~S.~Cook, S.~M.~Fleischer, and
B.~R.~Heckel,
``New Test of the Gravitational $1/r^2$ Law at Separations down to
52~$\mu$m,''
\emph{Phys.\ Rev.\ Lett.} \textbf{124}, 101101 (2020),
arXiv:2002.11761.

\bibitem{Ichiki}
K.~Ichiki, M.~Yahiro, T.~Kajino, M.~Orito, and G.~J.~Mathews,
``Observational Constraints on Dark Radiation in Brane Cosmology,''
\emph{Phys.\ Rev.\ D} \textbf{66}, 043521 (2002),
arXiv:astro-ph/0203272.

\bibitem{PerezVictoria}
M.~P\'erez-Victoria,
``Randall--Sundrum Models and the Regularized AdS/CFT Correspondence,''
\emph{JHEP} \textbf{05}, 064 (2001),
arXiv:hep-th/0105048.

\bibitem{AharonovBohm}
Y.~Aharonov and D.~Bohm,
``Significance of Electromagnetic Potentials in the Quantum Theory,''
\emph{Phys.\ Rev.} \textbf{115}, 485 (1959),
doi:10.1103/PhysRev.115.485.

\bibitem{PinheiroPires}
C.~Pinheiro and G.~O.~Pires,
``Extending the Barnes--Rivers Operators to $D=3$ Topological Gravity,''
\emph{Phys.\ Lett.\ B} \textbf{301}, 339 (1993),
arXiv:hep-th/9212008.

\bibitem{Paper1}
J.~Rodgers,
``Source--Readout Non-Equivalence: Descent and Equivariant Reconstruction Obstructions,''
Zenodo (2026),
\href{https://doi.org/10.5281/zenodo.21369967}{doi:10.5281/zenodo.21369967}.

\bibitem{Paper2}
J.~Rodgers,
``Target-Relative Necessity of Completion: When Readout Loss Obstructs, and What a Sufficient Extension Must Retain,''
Zenodo (2026),
\href{https://doi.org/10.5281/zenodo.21370501}{doi:10.5281/zenodo.21370501}.

\bibitem{Paper3}
J.~Rodgers,
``Canonical Minimal Source Completion: The Coarsest Readout Extension on which a Nominated Family of Source Relations Becomes Well Defined,''
Zenodo (2026),
\href{https://doi.org/10.5281/zenodo.21370763}{doi:10.5281/zenodo.21370763}.

\bibitem{Paper4}
J.~Rodgers,
``Geometric Realization of Completed Source Relations: Descent, Orbit Spaces, Invariant Relations, and Variational Response in Shadow Theory,''
Zenodo (2026),
\href{https://doi.org/10.5281/zenodo.21370985}{doi:10.5281/zenodo.21370985}.

\bibitem{Paper5}
J.~Rodgers,
``Observable Quotients and Exact Projected Dynamics: Closure, Memory, Minimal Dynamical Completion, and Effective Field Operators,''
Zenodo (2026),
\href{https://doi.org/10.5281/zenodo.21371251}{doi:10.5281/zenodo.21371251}.

\bibitem{Paper6}
J.~Rodgers,
``Non-Source Projection and Internal Identifiability,''
Zenodo (2026),
\href{https://doi.org/10.5281/zenodo.21371451}{doi:10.5281/zenodo.21371451}.

\end{thebibliography}

\end{document}
